\documentclass[11pt,letterpaper,reqno]{amsart}
% BEGIN INLINE manuscript-style
\makeatletter
\NeedsTeXFormat{LaTeX2e}[2020/10/01]
% Use with \documentclass[11pt,letterpaper,reqno]{amsart}.
% This file controls presentation only. It supplies no mathematical content.

\RequirePackage[T1]{fontenc}
\RequirePackage{lmodern}
\RequirePackage{microtype}
\RequirePackage[letterpaper,margin=1in,headheight=15pt]{geometry}
\RequirePackage{amsmath,amsthm,amssymb,mathtools}
\RequirePackage{mathrsfs}
\RequirePackage{enumitem}
\RequirePackage{perpage}
\RequirePackage{longtable}
\RequirePackage{xcolor}
\definecolor{manuscriptred}{rgb}{0.6,0.1,0.1}
\definecolor{manuscriptblue}{rgb}{0.1,0.3,0.5}
\definecolor{manuscriptmagenta}{HTML}{60445C}
\RequirePackage[breaklinks=true,colorlinks=true,
  linkcolor=manuscriptred,citecolor=manuscriptblue,urlcolor=manuscriptmagenta]{hyperref}
\RequirePackage{cleveref}
\hypersetup{pdfauthor={}}

% Retain the AMS paragraph, title, heading and running-head conventions.
% Equation counters restart each section: 2.i, 2.ii, ... .
% Appendix sections keep their letters: A.i, A.ii, ... .
\numberwithin{equation}{section}
\renewcommand{\theequation}{\thesection.\roman{equation}}
% Prose references include "eq." and omit parentheses around the number.
\renewcommand{\eqref}[1]{\textup{eq.~\ref{#1}}}
\crefname{equation}{eq.}{eqs.}
\Crefname{equation}{eq.}{eqs.}
\creflabelformat{equation}{#2#1#3}
\setcounter{tocdepth}{1}
\renewcommand{\thefigure}{\Roman{figure}}
\renewcommand{\thetable}{\Roman{table}}
\newcommand{\clneedspace}[1]{\par\begingroup
  \dimen0=#1\relax
  \vskip0pt plus\dimen0\penalty-100
  \vskip0pt plus-\dimen0\vskip\dimen0
  \penalty9999\vskip-\dimen0\vskip0pt\endgroup}
\MakePerPage{footnote}
\renewcommand{\thefootnote}{%
  \ifnum\value{footnote}>9\arabic{footnote}\else\fnsymbol{footnote}\fi}

% Upright statements. Main results are A, B, C; supporting statements
% share a section-local sequence. These end marks indicate typography,
% never proof completeness, validation, or a deferred-proof status.
\theoremstyle{definition}
\newtheorem{ms@theorem}{Theorem}
\renewcommand{\thems@theorem}{\Alph{ms@theorem}}
\newtheorem*{ms@theoremstar}{Theorem}
\newtheorem*{ms@theoremzero}{Theorem 0}
\newtheorem{ms@proposition}{Proposition}[section]
\newtheorem{ms@lemma}[ms@proposition]{Lemma}
\newtheorem{ms@definition}[ms@proposition]{Definition}
\newtheorem{ms@corollary}[ms@proposition]{Corollary}
\newtheorem*{ms@corollarystar}{Corollary}

\newcommand{\ms@majorstart}{%
  \pushQED{\qed}\renewcommand{\qedsymbol}{$\blacksquare\square$}}
\newcommand{\ms@minorstart}{%
  \pushQED{\qed}\renewcommand{\qedsymbol}{$\blacksquare$}}
\newenvironment{theorem}[1][]{\ms@majorstart\ms@theorem[#1]}
  {\popQED\endms@theorem}
\newenvironment{theorem*}[1][]{\ms@majorstart\ms@theoremstar[#1]}
  {\popQED\endms@theoremstar}
\newenvironment{theoremzero}[1][]{\ms@majorstart\ms@theoremzero[#1]}
  {\popQED\endms@theoremzero}
\newenvironment{proposition}[1][]{\ms@minorstart\ms@proposition[#1]}
  {\popQED\endms@proposition}
\newenvironment{lemma}[1][]{\ms@minorstart\ms@lemma[#1]}
  {\popQED\endms@lemma}
\newenvironment{definition}[1][]{\ms@minorstart\ms@definition[#1]}
  {\popQED\endms@definition}
\newenvironment{corollary}[1][]{\ms@majorstart\ms@corollary[#1]}
  {\popQED\endms@corollary}
\newenvironment{corollary*}[1][]{\ms@majorstart\ms@corollarystar[#1]}
  {\popQED\endms@corollarystar}
% The source's lemma/proposition variants are deliberately not carried over:
% choose one marker for each statement type rather than encode proof status.

\theoremstyle{remark}
\newtheorem{remark}[ms@proposition]{Remark}
\newtheorem*{remark*}{Remark}
% The ordinary amsthm proof environment retains its open square.

\crefname{ms@theorem}{theorem}{theorems}
\Crefname{ms@theorem}{Theorem}{Theorems}
\crefname{ms@proposition}{proposition}{propositions}
\Crefname{ms@proposition}{Proposition}{Propositions}
\crefname{ms@lemma}{lemma}{lemmas}
\Crefname{ms@lemma}{Lemma}{Lemmas}
\crefname{ms@definition}{definition}{definitions}
\Crefname{ms@definition}{Definition}{Definitions}
\crefname{ms@corollary}{corollary}{corollaries}
\Crefname{ms@corollary}{Corollary}{Corollaries}
\crefname{remark}{remark}{remarks}
\Crefname{remark}{Remark}{Remarks}

% Bibliography configuration for this manuscript.
\makeatother
% END INLINE manuscript-style

\title[Scalar continuum theory]{A scalar continuum theory from lattice Yang--Mills correlations}
\date{Working draft, 9 November 2026}
\begin{document}
\raggedbottom
\begin{abstract}
We assemble a scalar continuum hierarchy from uniform separated
moments of the Wilson action density. The lattice spacing tends
to zero along one admitted bare sequence after infinite-volume
Gibbs states have been chosen at each spacing. A quantitative
extension theorem controls the resulting moments on Schwartz
functions flat on all coincidence diagonals, with derivative
order linear in the number of fields and factorial coefficient
growth. Reflection positivity passes on the corresponding
positive-time test domain. Uniform lattice clustering estimates
give realizing pure states for a compact family of continuum
phases. The source insertion and physical two-point estimates
verify the stated general rotation and reconstruction criteria
for these same phases. Their finite means are retained, and the
connected short-distance lower bound yields a nonzero scalar
excitation. The actual uniform source, running and physical
comparison estimates remain hypotheses of the construction.
\end{abstract}
\maketitle
\tableofcontents

\section{Introduction}

The scalar Wilson action density can be tested at distinct
physical points without first defining products at coincident
points. Its continuum construction therefore begins with
separated smeared correlations. A common subsequence at those
correlations is only the first step: reflection positivity,
phase selection and reconstruction each use more than
pointwise existence of the individual limits.

The required quantitative information is a moment bound
uniform in the cutoff and the Gibbs states, with explicit
exponential and factorial growth in the number of insertions.
We use it first for simultaneous compactness.
A product-box Fourier construction then extends the moments
to Schwartz tests flat on the coincidence diagonals. Dyadic
separation cutoffs and Taylor flatness control the distribution
order and its dependence on the number of variables. The
positive-time version of the same approximation passes the
full reflected form, including its degree-zero terms.

The phase construction uses the all-state estimates before
taking continuum limits. The moment estimates control the
reflected sizes of local products. The lattice plateau law
then gives a physical clustering bound uniformly in the
cutoff. Evaluation coordinates provide both the
continuum phase law and actual realizing Gibbs-state sequences.
This realization allows the source insertion and connected
two-point estimates to be applied to each phase with the
original bare parameters and source normalization.

Ultraviolet stability under explicit running-coupling
assumptions is established in Ba{\l}aban's
\cite[Theorem~2 and Corollary~3, pp.~263--264]{Bal88}
and \cite[Theorem~1, pp.~355--356]{Bal89}.
The scalar construction additionally requires the
uniform source and physical comparison estimates
specified here.

The fixed-spacing inputs are supplied by
\cite{LP}, the trajectory by \cite{RT}, the graded source
and insertion estimates by \cite{SI}, and the physical
short-distance law by \cite{FP}. The general plateau and
angular continuation results are stated at their actual
interfaces below. We check the Euclidean test space,
growth and ordered positivity required by the corrected
Osterwalder--Schrader theorem \cite{OS75}.
The conclusion concerns the scalar sector generated by
the action density. Identifying a complete Yang--Mills
observable theory requires additional joint observable
constructions.

\section{Lattice states and the limiting test spaces}
\label{cl:sec-setting}

\subsection{The fixed-spacing state space}

Fix an odd integer \(L\ge3\). At spacing \(a>0\), choose one
orientation of each bond of \(a\mathbb Z^4\), and write
\(\Omega_a=SU(2)^{\mathcal B_a}\) with its product topology.
The reverse bond carries the inverse matrix. Haar measure has
total mass one. For each elementary square \(p\), let \(U_p\)
be its ordered boundary transport. Its starting vertex and
orientation do not affect
\[
 A_p=1-\tfrac12\operatorname{tr}U_p,\qquad 0\le A_p\le2.
\]
The trace is real for \(SU(2)\). We count each of the six
coordinate planes once at each lattice site and write \(x_p\)
for the plaquette barycentre.

For a finite bond set \(D\), the Wilson specification at
\(\beta\ge0\) holds all other bonds fixed and integrates the
bonds of \(D\) against normalized Haar measure with density
proportional to
\[
             \exp\!\left[-\beta
                \sum_{p:p\cap D\ne\varnothing}A_p\right].
\]
Here \(p\cap D\ne\varnothing\) means that the boundary of \(p\)
uses a bond of \(D\). A Wilson Gibbs state is a probability
measure satisfying these conditional expectation identities.
We denote the translation-invariant such states by
\(\mathcal G_{a,\beta}\). A periodic state is one invariant
under a finite-index lattice subgroup. The all-mirror theorem
of \cite[Theorem~B]{LP} implies that every periodic Wilson state
belongs to \(\mathcal G_{a,\beta}\), is invariant under signed
coordinate permutations, and is reflection positive at every
coordinate site, coordinate link and diagonal lattice mirror.
Its positivity applies to the full bounded Borel half-space
bond algebra.

\begin{proposition}[Fixed-spacing infinite-volume states]
\label{cl:prop-fixed-spacing}
For each fixed \(a>0\) and \(\beta\ge0\), periodic extensions
of cubic-torus Wilson measures have weak subsequential limits
as the torus side tends to infinity. Every such limit lies
in \(\mathcal G_{a,\beta}\). The set \(\mathcal G_{a,\beta}\)
is nonempty, compact and convex in the weak topology.
\end{proposition}
\begin{proof}
The countable product \(\Omega_a\) is compact and metrizable.
Its probability measures are therefore weakly sequentially
compact. The torus measures, extended periodically, are
invariant under every lattice translation. This invariance
passes to a weak limit because translations act continuously.

Fix \(D\) and a continuous cylinder function \(F\).
The conditional kernel \(\gamma_DF\) depends on only finitely
many bonds and is continuous. To see this, its denominator
is at least \(e^{-2\beta N(D)}\), where \(N(D)\) counts
plaquettes meeting \(D\); both numerator and denominator are
continuous parameter integrals over a compact group.
For sufficiently large tori, the bonds needed by \(D\),
its adjacent plaquettes and \(F\) embed without identifications.
The finite-torus conditional identity is then
\(\mu_\ell(\gamma_DF)=\mu_\ell(F)\). Weak convergence passes
this equality to the limit. Continuous cylinder functions
determine measures on the countable compact product, so the
identity extends to bounded Borel functions by the monotone
class theorem.

The equations for continuous cylinder functions are closed
conditions on probability measures. The translation equations
are closed too. Hence \(\mathcal G_{a,\beta}\) is closed
in a compact space. Each defining identity is linear in
the measure, which proves convexity.
\end{proof}

The proposition keeps \(a\) and \(\beta\) fixed. The following
cutoff assumptions supply the uniform estimates needed when
they vary.

\subsection{One bare sequence and its reference subtraction}

Take \(K_j\in\mathbb Z_{\ge0}\), \(a_j=L^{-K_j}\downarrow0\),
and nonnegative real bare couplings
\(\beta_j\). The data are assumed to lie on one admitted
trajectory family of \cite[Theorem~A]{RT}, with the same
fixed geometric and source constants. Assume in addition that the
ordinary source and actual rooted-domain hypotheses of
\cite[Theorems~B and C]{SI} hold uniformly on this family.
The short-distance statements will also use the actual
running, source-response and comparison assumptions of
\cite[Theorem~A and its cofinal alternative]{FP}. These
assumptions concern the given lattice theory and the
given bare sequence. They do not follow from compactness
of \(\mathcal G_{a,\beta}\).

Choose reference states \(\nu_j\in\mathcal G_{a_j,\beta_j}\).
At each datum, the integrating state
\(\rho_j\in\mathcal G_{a_j,\beta_j}\) may be different.
For \(f\in C_c^1(\mathbb R^4;\mathbb C)\), define
\begin{equation}\label{cl:eq-observable}
 \mathcal O_j(f)=\sum_p f(x_p)A_p,\qquad
 X_j^\nu(f)=\mathcal O_j(f)-\nu_j(\mathcal O_j(f)).
\end{equation}
The sum is finite. There is no extra power of \(a_j\).
This is the Wilson normalization used in \cite{SI,FP}.

A tuple \((f_1,\ldots,f_n)\) is separated if its nonzero
entries have pairwise disjoint compact supports. Such
supports have positive mutual distance. Set
\begin{equation}\label{cl:eq-cutoff-moments}
 G_{j,n}(f_1,\ldots,f_n)
   =\rho_j\!\left(\prod_{i=1}^n X_j^\nu(f_i)\right),
 \qquad G_{j,0}=1.
\end{equation}
We use complex multilinearity, not sesquilinearity, for
moments. Conjugation enters reflection-positive forms
separately. A tuple with a zero entry has value zero.
Write
\begin{equation}\label{cl:eq-sharp}
 \|f\|_\sharp=\max\{\|f\|_\infty,c_M\|\nabla f\|_\infty\},
 \qquad
 \|f\|_1=\max\{\|f\|_\infty,\|\nabla f\|_\infty\},
\end{equation}
where \(c_M>0\) is fixed.

The all-state graded estimate supplied by \cite[Theorem~B]{SI}
is the following. There are a fixed integer \(k\ge1\) and
\(B(R)=C(R+1)^4\) such that
\begin{equation}\label{cl:eq-graded-input}
 |G_{j,n}(f_1,\ldots,f_n)|
 \le B(R)^n(n!)^k\bar d^{-5n}
                         \prod_{i=1}^n\|f_i\|_\sharp .
\end{equation}
Each support is contained in its own radius-\(R\) ball;
the ball centres are arbitrary. For \(n\ge2\) the
supports have distance at least \(d>0\), and
\(\bar d=\min(d,1)\). For \(n=1\), set \(\bar d=1\).
The constants are independent of the cutoff, the two
states, the support positions and \(n\). They include
complex tests after a fixed enlargement of \(C\).
The same estimate applies when \(\rho_j=\nu_j\), and
when both are replaced by any other pair of states at
the same datum. This uniformity will permit phase
selection and self-centering without introducing a
new continuum sequence.

For each use of a varying constant \(C\), its allowed
dependencies are \(L,c_M\), the fixed source and
trajectory bounds, and explicitly displayed geometric
parameters. Such a constant does not depend on a
cutoff, state or field number. Dependence on the
field number is displayed through exponential and
factorial factors.

\clneedspace{23\baselineskip}
\subsection{Flat Schwartz functions}

\begin{definition}[The Euclidean test domains]
\label{cl:def-flat-spaces}
For \(n\ge2\), let
\(\Delta_n=\bigcup_{1\le i<h\le n}\{x_i=x_h\}\)
in \((\mathbb R^4)^n\), and put
\[
 \mathcal S_{\ne,n}
 =\{\varphi\in\mathcal S(\mathbb R^{4n}):
       \partial^\alpha\varphi|_{\Delta_n}=0
       \text{ for every multi-index }\alpha\}.
\]
Set \(\mathcal S_{\ne,1}=\mathcal S(\mathbb R^4)\)
and \(\mathcal S_{\ne,0}=\mathbb C\). These spaces
carry their Schwartz subspace topology. For \(r\)
a nonnegative integer and \(n\ge1\), use the seminorm
\begin{equation}\label{cl:eq-schwartz}
 |\varphi|_r
  =\max_{|\alpha|\le r}\sup_{x\in\mathbb R^{4n}}
       (1+|x|^2)^{r/2}|\partial^\alpha\varphi(x)|.
\end{equation}
For \(n=0\), the seminorm is the absolute value.
Let \(\mathcal S_{+,n}\subset\mathcal S_{\ne,n}\)
consist of functions whose every derivative vanishes
whenever at least one time coordinate is nonpositive.
Set \(\mathcal S_{+,0}=\mathbb C\). We use the fourth
coordinate as time: \(x=(\mathbf x,t)\), and
\(\theta(\mathbf x,t)=(\mathbf x,-t)\). Thus
\((\Theta\varphi)(x_1,\ldots,x_n)
 =\varphi(\theta x_1,\ldots,\theta x_n)\).
\end{definition}

These are closed subspaces: every pointwise derivative
evaluation is a continuous functional on Schwartz
space, and the defining conditions are intersections
of their kernels. Tensor products of compactly supported
smooth tests with pairwise disjoint supports belong to
\(\mathcal S_{\ne,n}\). Arbitrary tensor products need
not belong to it. In particular, continuity on this
domain assigns no value to a contact term supported
on a coincidence diagonal.

The two components of the construction have different
roles. The bound \eqref{cl:eq-graded-input} gives
compactness at every fixed separation and controls
the quantitative extension to the flat spaces.
The one-field insertion estimate of \cite[Theorem~C]{SI}
is an operator estimate uniform in the number of
surrounding fields; it is required separately by
the angular continuation criterion.


\section{A common subsequence for separated tests}
\label{cl:sec-compactness}

\begin{proposition}[Separated compactness]\label{cl:prop-separated-compactness}
Let \(a_j\to0\) be the given cutoff sequence. Suppose the moments
in \eqref{cl:eq-cutoff-moments} obey \eqref{cl:eq-graded-input} for
the integrating and centering states \((\rho_j,\nu_j)\), and for
\((\rho_j,\rho_j)\), with common constants. Assume both states
have the lattice translation and signed coordinate-permutation
symmetries of the fixed-spacing construction.
There is one subsequence on which every moment converges on every
tuple of separated compact \(C^1\) tests, at all orders.
Write \(G_n\) for these reference-centered limits and \(G_0=1\).
They obey the original bound, are multilinear on products of
fixed separated support spaces, and are symmetric, real and
invariant under all real translations and signed coordinate
permutations. Moreover \(G_1(f)=b\int f\) for a real constant \(b\).
The self-centered limits \(G_n^0\) and all connected limits of
order \(n\ge2\) exist on the same subsequence.

For any fixed product of separated compact support neighborhoods,
convergence is uniform on its \(C^1\)-compact sets of tuples.
It also holds for moving tests converging in \(C^1\) with supports
in such fixed neighborhoods. Closed inequalities involving
finitely many of these moment and connected evaluations pass
to the limit, with the same constants when these are uniform.
\end{proposition}

\begin{proof}
Consider bounded finite unions of open rational boxes in
\(\mathbb R^4\). They form a countable class. Separated compact
supports have neighborhoods \(U_i\) in this class whose closures
remain pairwise separated, by a finite cover of each support by
boxes of sufficiently small diameter. For each such \(U\), choose
smooth cutoffs \(\chi_m\) supported in \(U\) and equal to one on
the sets \(\{x\in U:\operatorname{dist}(x,U^c)\ge1/m\}\).
For example, convolve the indicator of the set at distance
greater than \(1/(2m)\) with a smooth nonnegative mollifier of
mass one and support radius \(1/(8m)\).

Use the rational complex spans of \(\chi_m\) times the real sines
and cosines of integral frequencies on integer-period boxes.
The resulting family is countable and dense in \(C_c^1(U)\)
for the \(C^1\) norm. Indeed, for a given test choose \(\chi_m=1\)
on a neighborhood of its support. Mollify the test on a smaller
scale and periodize it on a box containing \(\overline U\) and
the mollified support in its interior. Uniform continuity of the
first derivatives gives \(C^1\) convergence of the mollification.
Integration by parts in the Fourier series of the smooth function
gives absolute summability after one differentiation. Approximate
the coefficients of a finite partial sum by rational complex
numbers and multiply by the fixed cutoff. Multiplication by it
is bounded in \(C^1\) and leaves the target test unchanged.

List all finite tuples from these countable families with pairwise
separated neighborhood closures, including the one-point tuples.
Their values are bounded by \eqref{cl:eq-graded-input} at their
fixed order, radius and separation. Successive subsequences
converge at the successive entries of this countable list.
Selecting an increasing index from the \(r\)-th nested subsequence
at stage \(r\) gives one sequence converging on the entire list.
An eventual bound on each fixed neighborhood product suffices
for this selection.

To extend convergence, let \(\mathbf f,\mathbf h\) lie in the same
product of fixed separated support neighborhoods. Telescope by
replacing the entries one at a time. Every intermediate tuple
remains separated. Thus for the fixed order \(n\)
\begin{equation}\label{cl:eq-moment-telescoping}
 \bigl|G_{j,n}(\mathbf f)-G_{j,n}(\mathbf h)\bigr|
 \le C\sum_{i=1}^n \|f_i-h_i\|_\sharp
      \prod_{l\ne i}\max\{\|f_l\|_\sharp,\|h_l\|_\sharp\},
\end{equation}
where \(C\) depends on \(n\) and the radius and separation of the
neighborhoods. Approximate a given tuple in its own neighborhoods
by a listed tuple. Inequality~\eqref{cl:eq-moment-telescoping}
makes this approximation uniform in the cutoff tail. The values
of the original tuple are Cauchy along the same subsequence.
The same inequality proves uniform convergence on compact sets
by finite \(C^1\) nets, and convergence on moving tuples by
comparison with their limit. It also proves continuity of \(G_n\)
on each fixed support product. The graded bound passes on the
original supports, so its constants use their radius and gap.

Multilinearity on these products, symmetry and reality pass from
the lattice identities. For a real translation \(v\), choose
\(v_j\in a_j\mathbb Z^4\) with \(v_j\to v\). For
\((\tau_v f)(x)=f(x-v)\) one has
\[
 \|\tau_{v_j}f-\tau_v f\|_\sharp
 \le\max\{\|\nabla f\|_\infty |v_j-v|,
                  c_M\omega_{\nabla f}(|v_j-v|)\},
\]
where \(\omega_{\nabla f}\) is the uniform modulus of continuity
of \(\nabla f\). Fixed separated neighborhoods of the limiting
translated supports contain the moving supports for all large
\(j\). Lattice translation invariance and moving-test convergence
give invariance under \(v\). The signed coordinate permutations
act exactly at each cutoff. The same moving-test conclusion
applies to pullbacks by converging Euclidean isometries whenever
the limiting tuple remains separated, by the chain rule and
uniform continuity.

At one point the difference of the two plaquette means is the
same number \(\delta_j\) for all positions and planes. This uses
the signed coordinate permutations as well as translations.
Put
\begin{equation}\label{cl:eq-reference-mean}
 b_j=6a_j^{-4}\delta_j,\qquad
 R_j(f)=\frac{a_j^4}{6}\sum_p f(x_p),\qquad
 G_{j,1}(f)=b_jR_j(f).
\end{equation}
For each of the six plane labels the shifted grid gives a Riemann
sum for \(\int f\). The mesh cubes meeting a support of radius
\(R\) have total volume at most \(C(R+2a_j)^4\).
The oscillation in one cube is at most
\(2a_j\|\nabla f\|_\infty\), so
\[
 |R_j(f)-\textstyle\int f|
 \le C(R+2a_j)^4a_j\|\nabla f\|_\infty.
\]
Choose a fixed compact test of integral one and radius at most
one. Its Riemann sum is at least \(1/2\) eventually, independently
of the two states. The one-point bound then bounds \(b_j\)
uniformly over the states. Convergence at this test gives
\(b_j\to b\in\mathbb R\) along the chosen subsequence.
The same identity then gives \(G_1(f)=b\int f\) for every test.

For a finite index set write \(\mathbf f_A=(f_i)_{i\in A}\) in
its inherited order. Self-centering is obtained by subtracting
\(G_{j,1}(f_i)\) from each reference-centered variable. Expanding
the product and passing to the limit gives
\begin{equation}\label{cl:eq-self-centered-moments}
 G_n^0(\mathbf f)
  =\sum_{A\subseteq\{1,\ldots,n\}}G_{|A|}(\mathbf f_A)
                \prod_{i\notin A}\left(-b\int f_i\right).
\end{equation}
Its graded bound passes directly from the bound for
\((\rho_j,\rho_j)\). In particular \(G_0^0=1\) and \(G_1^0=0\).
For \(n\ge2\) the cumulant is unchanged by deterministic shifts.
Its limit is the finite partition sum
\begin{equation}\label{cl:eq-separated-connected}
 S_n^T(\mathbf f)=
 \sum_{\pi}(-1)^{|\pi|-1}(|\pi|-1)!
        \prod_{A\in\pi}G_{|A|}(\mathbf f_A).
\end{equation}
Here \(\pi\) ranges over partitions of \(\{1,\ldots,n\}\).
The formula follows from expanding the logarithm of the normalized
source generating function at zero, or equivalently the
moment--cumulant identity of \cite[Lemma~2.1]{SI}.
In particular \(G_n^0\) is the sum of products of the \(S_l^T\)
over partitions with all block sizes at least two.

The same factorial exponent works for the connected functions.
Indeed we may take the exponent \(k\) of
\eqref{cl:eq-graded-input} to be at least one. The moment
weights of a partition into sizes \(n_1,\ldots,n_r\) carry
\(\prod_l(n_l!)^k\le (n!)^{k-1}\prod_l n_l!\).
Summing \(\prod_l n_l!\) over ordered blocks of any one composition
of \(n\) gives \(n!\), since the number of assignments is
\(n!/\prod_l n_l!\). Each unordered partition with \(r\) blocks
occurs \(r!\) times. The sum weighted by \((r-1)!\) is therefore
\(n!\binom{n-1}{r-1}/r\). Summing over \(r\) gives at most
\(2^{n-1}(n!)^k\) after the common \((n!)^{k-1}\) factor.
The remaining moment bounds multiply to the original product
norm and separation power. Hence one extra factor \(2^n\) in
the exponential constant suffices.

All the limits used here are simultaneous on each fixed finite
collection of tuples and their subtuples. Thus inequalities
defined by closed conditions on such evaluations, with any
convergent scalar coefficients, pass by continuity. This also
applies to moving tuples by \eqref{cl:eq-moment-telescoping}.
\end{proof}

\begin{corollary}[Torus sequences]\label{cl:cor-torus-compactness}
Consider compatible cubic Wilson tori with \(a_j\to0\) and
even bare lattice side counts \(\ell_j/a_j\).
Use their Wilson expectation for both integration and centering.
On a fixed torus, uniform bounds of the form
\(C_{n,R,d}\prod_i\|f_i\|_\sharp\) at each order and support
radius and separation suffice for one subsequence on all
separated \(C^1\) tuples. The limiting hierarchy is self-centered
and has all torus translations and signed coordinate-permutation
symmetries. If the torus sides tend to infinity, a fixed compact
tuple is transported in a common injective lift. Under the
finite-torus graded bounds with uniform constants, the same
construction gives the separated limits and bounds of
Proposition~\ref{cl:prop-separated-compactness} on \(\mathbb R^4\).

Suppose instead that a self-centered continuum hierarchy has
already been extracted at each of a sequence of torus sides
tending to infinity. A separated limit of those hierarchies,
with the uniform per-support bounds at finite cutoffs, is also
realized by one sequence of original torus data with meshes
tending to zero.
\end{corollary}

\begin{proof}
On a fixed torus use a countable basis of rational coordinate
boxes in the period-normalized coordinates, the same
support-adapted cutoffs and rational Fourier polynomials.
The preceding extraction and telescoping proof applies there.
Translations are rounded modulo the period. The one-point
functions vanish because the tori use self-centering.
On growing tori any fixed finite family of compact support
neighborhoods fits, eventually, in a common lift of diameter
less than a quarter of the side. Periodizing the tests in this
lift preserves the support distances and \(C^1\) norms.
Indeed the copies are disjoint, and the quotient is an isometry
on the chosen lift. Thus the bounds needed for each entry of
the countable list hold eventually. The moving-test and
symmetry arguments then apply on the lifted support
neighborhoods with their fixed margins.

For the last assertion fix the countable list of separated
Euclidean tuples used above. On the \(j\)-th torus select one
term of its defining cutoff subsequence with \(K\ge j\) whose
values differ by at most \(1/j\) from the fixed-torus limits
on the first \(j\) listed tuples whose neighborhoods embed
there. There are finitely many such conditions, and each
holds eventually along the defining subsequence at that torus.
Every fixed listed tuple eventually embeds and is among those
tested. The selected torus values and the iterated limits
therefore agree on this list. Inequality~\eqref{cl:eq-moment-telescoping}
with the common local constants applies both to the inner
limits and the selected finite cutoffs. Approximate any fixed
tuple in its separated neighborhoods, and let this approximation
error tend to zero after the diagonal limit. This gives agreement
on every fixed separated tuple and all its subtuples.
\end{proof}

\clneedspace{22\baselineskip}
\subsection{Recognition on sufficiently large tori}

\begin{proposition}[Large torus recognition]
\label{cl:prop-large-torus-recognition}
For each fixed bare datum \(K\in\mathbb Z_{\ge0}\), \(\beta\ge0\),
with \(a_K=L^{-K}\),
fix one compatible cofinal family of Wilson tori and let
\(\mathcal G^{\rm tor}_{K,\beta}\)
be the set of weak limit points of their periodic lifts.
Write \(\mu_{\rm torus}\) for the periodic lift at the
compared side.
There is a side threshold \(\ell_\flat(K,\beta)\ge2L^K\)
with the following properties. For every side in that family
above the threshold there is a state
\(\mu\in\mathcal G^{\rm tor}_{K,\beta}\) whose marginal on the
bonds of the physical max-norm cube of radius \(K+1\) has
Haar density ratio to the torus marginal within
\(\epsilon_K=2e^{-1/a_K}\) of one uniformly.

For \(K\to\infty\) on such sides, replacing each torus law by
its chosen state leaves the limiting connected correlations
and self-centered moments of every fixed separated compact
tuple unchanged. Conversely, a sequence of states in these
torus limit sets at their respective bare data admits torus
approximants above the same thresholds with the same limits.
On the chosen cube the torus diagonal reflection form satisfies
\begin{equation}\label{cl:eq-recognition-rp}
 \mu_{\rm torus}(\Theta_\Pi\overline F\,F)
 \ge -\epsilon_K\,\mu_{\rm torus}(|F|^2)
\end{equation}
for a bounded Borel \(F\) on one side of a diagonal mirror
\(\Pi\) through the origin. Here \(\Theta_\Pi\) is the linear
pullback by the bond reflection; conjugation is written separately.
\end{proposition}

\begin{proof}
At fixed spacing first use any finite bond set \(W\) such that
it and the exterior bonds of its incident plaquettes embed in
the torus. Write \(h_W(V\mid U)\) for the conditional density
in the Wilson specification of
Proposition~\ref{cl:prop-fixed-spacing}, relative to normalized
Haar on \(W\). It is continuous in the finitely many interior
and exterior variables. If \(N_W\) is the number of plaquettes
meeting \(W\), then \(h_W\ge e^{-2\beta N_W}\), since every
plaquette action lies in \([0,2]\).
For the periodic lift \(\widetilde\mu_\ell\) and each limit
state \(\mu\) the respective marginals have densities
\[
 m_{\widetilde\mu_\ell}(V)=\int h_W(V\mid U)\,d\widetilde\mu_\ell(U),
 \qquad
 m_\mu(V)=\int h_W(V\mid U)\,d\mu(U).
\]
The first formula follows by integrating the embedded window
in the finite Wilson integral. The second is the DLR identity.
The family \(h_W(V\mid\cdot)\) has compact image in the sup norm
as \(V\) ranges over a compact product of copies of \(SU(2)\).
Thus along any subsequence of torus lifts converging weakly to
\(\mu\), a finite sup-norm net of this family gives uniform
convergence of their marginal densities to \(m_\mu\).
The common positive lower bound gives uniform convergence
of \(m_\mu/m_{\widetilde\mu_\ell}\) to one on that subsequence.
It follows that
\[
 \inf_{\mu\in\mathcal G^{\rm tor}_{K,\beta}}
 \|m_\mu/m_{\widetilde\mu_\ell}-1\|_\infty\longrightarrow0
 \qquad(\ell\to\infty).
\]
Otherwise a violating sequence of sides would have a weak
subsequence with a limit in this same set.

Now take all bonds with both endpoints in the radius-\((K+1)\)
cube, with either orientation represented once.
Choose \(\ell_\flat\) so that the window with its plaquette
collar embeds and the last infimum is at most \(e^{-1/a_K}\)
on every later side. Increasing it to \(2L^K\) if needed gives
the stipulated growth. Select one state within \(2e^{-1/a_K}\)
in the density-ratio norm. For every bounded measurable \(B\)
of the window,
\begin{equation}\label{cl:eq-recognition-density}
 |\mu(B)-\widetilde\mu_\ell(B)|
 \le\epsilon_K\,\widetilde\mu_\ell(|B|).
\end{equation}
This is one window estimate for all the observables used below.

Fix finitely many compact tests in a cube of physical radius
\(R\), with a plaquette margin inside the window for large \(K\).
The number of plaquette centres per test is at most
\(C(1+R/a_K)^4\). Hence \eqref{cl:eq-recognition-density}
bounds a difference of raw moments of order \(n\) by
\[
 C^n\epsilon_K(1+R/a_K)^{4n}\prod_i\|f_i\|_\infty .
\]
Expansion over subsets for self-centering and over partitions
for cumulants gives the same estimate with a constant
depending on the fixed order. For a product of expectations,
replace its factors one at a time and use the corresponding
raw bounds for the other factors. The exponential
\(e^{-1/a_K}\) dominates every fixed power of \(a_K^{-1}\).
This proves recognition of all fixed-tuple limits, with the
same chosen state for every tuple eventually contained in the
window. For the converse, at each fixed datum use a cofinal
subsequence converging to its prescribed torus limit point.
The ratio convergence already proved allows a side as large
as desired at that datum with the same window tolerance.
The same moment estimate gives the converse.

For the reflection claim, the cube and its bond set are
invariant under a diagonal reflection through the origin.
The torus lift is invariant under this signed permutation,
so its expectation of \(\Theta_\Pi\overline F\,F\) is real.
The state \(\mu\) is reflection positive by
\cite[Theorem B]{LP}. Apply \eqref{cl:eq-recognition-density}
to \(\Theta_\Pi\overline F\,F\) and use
\[
 \widetilde\mu_\ell(|\Theta_\Pi\overline F|\,|F|)
 \le \widetilde\mu_\ell(|F|^2)
\]
by Cauchy--Schwarz and reflection invariance. This gives
\eqref{cl:eq-recognition-rp}. Its error tends to zero for
polynomials of fixed degree in fixed compactly supported
observables and their centered versions, by the same
plaquette count. For an affine lattice diagonal mirror,
a lattice translation moves it to the origin. On fixed
bounded support sets and bounded mirror offsets the translated
observables lie in the same growing window eventually.
\end{proof}

\section{Quantitative extension to flat Schwartz tests}
\label{cl:sec-flat-extension}

Use the seminorm \eqref{cl:eq-schwartz} and the spaces of
Definition~\ref{cl:def-flat-spaces}, with the one-slot norm
\(\|f\|_1\) fixed in \eqref{cl:eq-sharp}.

\begin{proposition}[Flat extension and separated tensor density]
\label{cl:prop-flat-extension}
Let \(F_n\), \(n\ge1\), be forms on separated compact \(C^1\) tuples,
multilinear on each product of separated support spaces.
Suppose \(D\ge1\), \(\kappa\ge0\) and \(\alpha\ge0\) are independent
of \(n\), and for radius-one supports at mutual distance at least
\(d>0\) one has
\begin{equation}\label{cl:eq-flat-tensor-input}
 |F_n(f_1,\ldots,f_n)|\le
 D^n(n!)^\kappa\min(1,d)^{-\alpha n}\prod_i\|f_i\|_1.
\end{equation}
At \(n=1\) take \(d=1\). Set \(A=\lceil\alpha\rceil\) and
\(k_*=A+15\). There is a unique continuous linear extension to
\(\mathcal S_{\ne,n}\) agreeing with the values on smooth compact
separated tensors, and
\begin{equation}\label{cl:eq-flat-growth}
 |F_n(\varphi)|\le C^n(n!)^{\kappa+32}|\varphi|_{k_*n},
 \qquad n\ge1 .
\end{equation}
Here \(C\) depends on \(D,\alpha\) and numerical cutoff choices,
uniformly in \(n\) and all families obeying
\eqref{cl:eq-flat-tensor-input} with those constants.
Finite sums of the indicated tensors are dense in
\(\mathcal S_{\ne,n}\) in the Schwartz topology.
For the positive-time flat subspace, their entries may all have
compact support strictly in the open positive half-space.
This density assertion also holds for any affine half-space
by an isometric change of test coordinates.
The radius-one bound also gives a radius-\(R\) bound in
\eqref{cl:eq-flat-tensor-input} with \(CD(1+R)^4\) in place of
\(D\).
\end{proposition}

\begin{proof}
\emph{Cutoffs.}
We use a finite set of fixed smooth cutoff functions valued in
\([0,1]\) whose derivatives of order \(q\) are bounded by
\(C^{q+1}(q!)^2\). To construct them, start with
\(a(t)=e^{-1/t}\) for \(t>0\) and zero for \(t\le0\).
On the complex circle \(|z-t|=t/2\) one has
\(\operatorname{Re}(1/z)\ge2/(9t)\). Cauchy's formula gives
\[
 |a^{(q)}(t)|\le q!(2/t)^q e^{-2/(9t)}.
\]
The maximum of \(t^{-q}e^{-c/t}\) and the inequality
\((q/e)^q\le q!\) give the stated bounds and flatness at zero.
The smooth step \(s(t)=a(t)/(a(t)+a(1-t))\) is zero at \(t\le0\)
and one at \(t\ge1\). On \([0,1]\) its denominator is at least
\(e^{-2}\). A reciprocal of a denominator bounded below by
\(\delta>0\) has the same derivative bound. Indeed, differentiating
\(D D^{-1}=1\) and dividing by \((q!)^2\) introduces the
coefficients \(1/\binom qr\). If the denominator bound uses a
constant \(C_D\), the induction for the reciprocal closes with
\(M B^q(q!)^2\), \(M\ge\delta^{-1}\), as soon as
\((C_D/\delta)\sum_{r\ge1}(C_D/B)^r\le1\).
Leibniz's rule then gives the bound for the step.

Take an even cutoff \(\eta(t)\) equal to one for \(|t|\le5/4\)
and zero for \(|t|\ge2\), using \(s((2-|t|)/(3/4))\) on the
transition.
Put \(\psi(t)=\eta(2t)\) and
\(\chi_1(t)=\psi(t)/\sum_{m\in\mathbb Z}\psi(t-m)\).
The denominator is at least one by the nearest integer, and its
locally finite sum uses at most three terms and has the same
derivative bounds. Define on \(\mathbb R^4\)
\[
 \chi(u)=\prod_{\mu=1}^4\chi_1(u_\mu),\qquad
 \chi^+(u)=\prod_{\mu=1}^4\eta(u_\mu),\qquad
 \zeta(u)=1-\chi^+(u).
\]
The integer translates of \(\chi\) sum to one.
Its support is in \([-1,1]^4\), where \(\chi^+=1\), and
\(\chi^+\) has support in \([-2,2]^4\). Thus \(\zeta\) is zero
on the first cube and one outside the second. All these
cutoffs have the stated derivative bounds. For \(N\) such
factors at a scale proportional to \(h\), Leibniz's rule yields
\begin{equation}\label{cl:eq-cutoff-product}
 \left|\partial^\beta \prod_{\nu=1}^N b_\nu\right|
 \le (CN)^q(q!)^2h^{-q},\qquad q=|\beta|.
\end{equation}
To see the dependence on \(N\), distribute the derivative
among the factors. The multinomial coefficients sum to \(N^q\).
The assigned factorials have product at most \(q!\), and
at most \(q\) factors are differentiated. Each undifferentiated
factor is bounded by one. Using a difference \(x_i-x_l\)
as an argument changes only derivative signs in this estimate.

For the radius-\(R\) assertion, use translates of this partition
at a fixed mesh small enough for each piece to fit in a unit
ball. At most \(C(1+R)^4\) pieces meet a support of radius \(R\).
Their \(C^1\) norms are bounded by a fixed multiple of the original
norm. Multiplication by a cutoff preserves containment in the
original support, and thus preserves each interslot gap.
Multilinearity and the radius-one estimate give the assertion.

\emph{Local product boxes.}
Let \(0<h\le1/16\) and \(\lambda=h/64\).
Suppose the cubes \(c_i+[-2\lambda,2\lambda]^4\) have mutual
distance at least \(h/2\). For \(u\) smooth and compactly
supported in the product of open inner cubes
\(c_i+(-5\lambda/4,5\lambda/4)^4\),
expand it on the product of the period-\(8\lambda\) cubes.
Multiply each one-slot Fourier exponential by
\(\chi^+((x_i-c_i)/\lambda)\). The resulting product series
equals \(u\), since the periodic extension agrees with \(u\)
on the support of the product cutoff.
For its Fourier coefficients \(\widehat u(k)\),
\(k=(k_i)_{i=1}^n\), put
\[
 J_{r,\lambda}(u)=
     \max_{|\beta|\le r}\lambda^{|\beta|}
                   \|\partial^\beta u\|_\infty .
\]
Rescale each slot to a period-two cube. Integration by parts with
\((1-\pi^{-2}\Delta_i)^4\) in each slot bounds the coefficients by
\[
 |\widehat u(k)|\le C^n J_{8n,\lambda}(u)
                    \prod_i(1+|k_i|^2)^{-4}.
\]
The coefficient sum of this differential operator is a fixed
constant per slot. The Fourier normalization is the inverse
volume of the period cube, so it does not add a power of
\(\lambda^{-1}\). Each localized exponential has
\(\|\cdot\|_1\le C\lambda^{-1}(1+|k_i|)\).
The single-slot sum of \((1+|k_i|)(1+|k_i|^2)^{-4}\)
over \(\mathbb Z^4\) is finite. Equation~\ref{cl:eq-flat-tensor-input} therefore defines a local
functional \(T(u)\) by an absolutely convergent series and gives
\begin{equation}\label{cl:eq-flat-box-bound}
 |T(u)|\le (CD)^n(n!)^\kappa
                  h^{-(\alpha+1)n}J_{8n,\lambda}(u).
\end{equation}

For a tensor in the inner cubes the Fourier coefficients factor.
The localized one-slot series converge in \(C^1\) on fixed
separated support spaces. The multilinear bound shows that
the local functional takes the prescribed tensor value.
On overlaps, partition a compactly supported test into smaller
product boxes whose expanded supports are inside the overlap.
This is possible since the overlap is open and disjoint from
the diagonals. In each smaller box the same cutoff Fourier
series converges with all derivatives on a fixed compact set,
by integration by parts at any fixed order. Both local
functionals agree on these tensor approximants and are continuous
by \eqref{cl:eq-flat-box-bound}, so they agree on the test.
Finite smooth partitions of unity now define a single functional
on compact tests off the diagonals. A different partition gives
the same value by refining both partitions and using agreement
on overlaps. For an arbitrary original separated smooth tensor,
partition each slot into smaller compact pieces with their
expanded boxes still separated across slots. Multilinearity on
the original support product then gives the prescribed value.
In particular linear relations among separated tensors have
zero value under this construction.

\emph{Summing over scales and positions.}
Let \(h_j=2^{-j}/16\), \(\lambda_j=h_j/64\), and set
\[
 P_h(x)=\prod_{i<l}\zeta((x_i-x_l)/h),\qquad
 D_0=P_{h_0},\qquad D_j=P_{h_j}-P_{h_{j-1}}\quad(j\ge1).
\]
Off the diagonals the sum of the \(D_j\) is one. A fine layer
and all its derivatives are supported where the smallest
max-norm pair distance lies between \(h_j\) and \(4h_j\).
At \(n=1\) the pair product is one, so there are no fine layers.
At level \(j\) use the product partition
\(\chi_{j,z}(x)=\prod_i\chi(x_i/\lambda_j-z_i)\) for
\(z=(z_i)_{i=1}^n\in(\mathbb Z^4)^n\).
When \(\chi_{j,z}D_j\) is nonzero, a point in its support
has every pair distance at least \(h_j\). The slot centres
are at pair distance at least \(h_j-2\lambda_j\) in max norm.
Their expanded cubes are separated by at least
\(h_j-6\lambda_j>h_j/2\).
The closed grid supports lie strictly inside the open inner
cubes used in the local construction.
Thus \eqref{cl:eq-flat-box-bound} applies to
\(u=\chi_{j,z}D_j\varphi\).

In each layer there are at most \(C n^2\) cutoff factors.
For derivatives through order \(8n\), \eqref{cl:eq-cutoff-product}
and the scale factor in \(J_{8n,\lambda_j}\) give a cutoff cost
bounded by
\begin{equation}\label{cl:eq-factorial-accounting}
 (Cn^2)^{8n}((8n)!)^2 \le C^n(n!)^{32}.
\end{equation}
Here \(n^{16n}\le e^{16n}(n!)^{16}\), and
\((8n)!\le8^{8n}(n!)^8\) by the multinomial formula.
No dimension-dependent constant is used in place of this bound.

On a fine layer move the first point of a closest pair onto
the second. The resulting point is on a diagonal at Euclidean
distance at most \(8h_j\). The segment uses at most four
coordinate directions. Since all derivatives of \(\varphi\)
vanish at its diagonal endpoint, Taylor's formula gives, for
\(|\beta|\le8n\) and \(m=(A+5)n+2\),
\[
 |\partial^\beta\varphi(x)|\le
 C_A^n h_j^m (1+|x|^2)^{-k_*n/2}|\varphi|_{k_*n}.
\]
Indeed \(8n+m\le k_*n\). The directional remainder has
coefficient at most \((Ch_j)^m/m!\), using the four-coordinate
\(\ell^1\) displacement. The weights along a segment of length
at most \(1/2\) are comparable to the weight at \(x\), at cost
\(C^{k_*n}\). Both costs are bounded by \(C_A^n\) after
retaining the displayed power of \(h_j\).
For \(j=0\) use the direct seminorm bound and omit \(h_j^m\).

By Leibniz's rule, the scale factor in \(J_{8n,\lambda_j}\)
cancels the inverse powers of \(h_j\) on the cutoff derivatives.
Combining the last estimates gives
\[
 J_{8n,\lambda_j}(\chi_{j,z}D_j\varphi)
 \le C_A^n(n!)^{32}h_j^m|\varphi|_{k_*n}
     \sup_{x\in\operatorname{supp}\chi_{j,z}}
                 (1+|x|^2)^{-k_*n/2}
\]
on a fine layer, with the same coarse-layer convention.
For the grid sum use
\(1+|x|^2\ge\prod_i(1+|x_i|^2)^{1/n}\).
Each \(x_i\) is a bounded distance from its grid centre, so
\[
 \sum_z\sup_{x\in\operatorname{supp}\chi_{j,z}}
                 (1+|x|^2)^{-k_*n/2}
 \le C_A^n \prod_i\sum_{z_i\in\mathbb Z^4}
               (1+|\lambda_j z_i|^2)^{-k_*/2}
 \le C_A^n h_j^{-4n}.
\]
The one-slot bound follows by integrating over its mesh cubes,
since \(k_*>4\). Substitution in \eqref{cl:eq-flat-box-bound}
gives the fine-scale bound
\[
 \sum_z |T(\chi_{j,z}D_j\varphi)|
 \le C^n(n!)^{\kappa+32}|\varphi|_{k_*n}
                  h_j^{m-(\alpha+5)n}.
\]
The power of \(h_j\) is at least two. At \(j=0\) its fixed scale
cost is absorbed in \(C^n\). The sum in \(j\) is absolute,
so it defines the extension and proves \eqref{cl:eq-flat-growth}.
For compact support away from the diagonals this sum is finite
and is the glued local functional.

\emph{Density and the time boundary.}
First multiply by \(\prod_i\chi^+(x_i/R)\) and send \(R\to\infty\).
For a fixed seminorm, the discarded derivatives have a point
with \(|x_i|\ge R\), and one extra power of Schwartz decay
bounds their seminorm by \(C_{n,r}R^{-1}|\varphi|_{r+1}\).
Next multiply by \(P_h\) and send \(h\downarrow0\).
Every surviving derivative of \((1-P_h)\varphi\) has a pair
within max-norm distance \(2h\). For fixed \(n,r\), the preceding
Taylor argument with any \(M>r\) gives a bound
\(C_{n,r,M}h^{M-r}|\varphi|_{r+M}\).
Thus compact tests away from the diagonals are dense.
Apply a finite product-box partition and the cutoff Fourier
expansions to those tests. Partial sums converge with all
derivatives on a common compact set, hence in the Schwartz
topology. Choosing the cutoffs and finite sums successively
for the first \(r\) seminorms gives a single sequence of finite
sums of separated tensors.

On \(\mathcal S_{+,n}\) also multiply by
\(B_h(x)=\prod_i s(x_{i,4}/h-1)\).
The cutoff is zero when a time is at most \(h\) and one when
all are at least \(2h\). Where it removes a nonzero derivative
of a positive-time test, some time lies in \([0,2h]\).
Taylor expansion to that time hyperplane gives the same
decay \(h^{M-r}\), because every jet there vanishes.
The compact supports of the approximants lie strictly in
the positive half-space, so their enlarged Fourier boxes
can be chosen there too. If one uses \(P_hB_h\) directly
in the quantitative dyadic construction, a fine layer has
one close pair or one time within \(4h\) of zero.
The same Taylor and cutoff estimates apply, since there
are still at most \(Cn^2\) factors. The constants in
\eqref{cl:eq-flat-growth} retain the same order and factorial
exponent on that subspace. An affine isometry of the tests
gives the density statement for another half-space.

Continuity and this density give uniqueness, independently
of the boxes and cutoffs. They also transport any symmetry
of the separated forms whose test pullback is continuous
in the Schwartz topology.
\end{proof}

\begin{lemma}[Integration of omitted slots]\label{cl:lem-flat-marginal}
For \(A\subseteq\{1,\ldots,n\}\), let \(I_A\) integrate the
variables indexed by \(A\), leaving the other slots in order.
It maps \(\mathcal S_{\ne,n}\) continuously into
\(\mathcal S_{\ne,n-|A|}\) and preserves positive-time
flatness in the remaining variables. If \(m=|A|\) and \(r\ge0\)
is an integer, then
\begin{equation}\label{cl:eq-marginal-bound}
 |I_A\varphi|_r\le C_6^{\,m}|\varphi|_{r+6m},
 \qquad C_6=\int_{\mathbb R^4}(1+|y|^2)^{-3}\,dy.
\end{equation}
For no remaining variables the left side is the absolute value
of the scalar integral.
\end{lemma}

\begin{proof}
Let \(x\) collect the remaining variables and \(y_i\) the
integrated ones. The weight inequality
\[
 (1+|x|^2)^{r/2}\prod_i(1+|y_i|^2)^3
 \le(1+|x|^2+\textstyle\sum_i|y_i|^2)^{(r+6m)/2}
\]
gives an integrable majorant for every retained derivative of
order at most \(r\). Differentiation under the integral and
integration of the product weights yield \eqref{cl:eq-marginal-bound}.
On a remaining diagonal or a forbidden retained time, the
corresponding derivatives vanish at every integrated datum.
They therefore vanish after integration as well.
\end{proof}

\begin{corollary}[Growth and convergence of the extended hierarchy]
\label{cl:cor-flat-growth}
The separated limits of Proposition~\ref{cl:prop-separated-compactness}
extend uniquely to the flat spaces. Their full and self-centered
moments, and their higher connected functions, satisfy
\begin{equation}\label{cl:eq-hierarchy-growth}
 |F_n(\varphi)|\le C^n(n!)^{k+32}|\varphi|_{20n},
 \qquad n\ge1 ,
\end{equation}
uniformly for limits from the same graded input constants.
For infinite-lattice sequences the same subsequence converges
on every flat Schwartz test after extending the cutoff forms.
For growing tori the conclusion concerns the extended limit
of the compact tuples. The fixed-torus limits define unique
distributions off the coincidence diagonals.
\end{corollary}

\begin{proof}
On radius-one supports \(\|f\|_\sharp\le\max(1,c_M)\|f\|_1\).
The graded input and the cumulant estimate therefore give
\eqref{cl:eq-flat-tensor-input} with \(\alpha=5,\kappa=k\)
and a common \(D\), after increasing it by a fixed factor.
Proposition~\ref{cl:prop-flat-extension} gives
\eqref{cl:eq-hierarchy-growth}. It applies with the same
constants to any admitted phase realized on a sequence of
states with the same two-state estimates, including its
reference-centered one-point function.

For the infinite-lattice sequence the uniform graded bound
also applies to each cutoff form. Approximate a fixed
\(\varphi\) by a finite sum of smooth separated tensors in
\(|\cdot|_{20n}\). On that finite sum the cutoff limits are
already known. Equation~\ref{cl:eq-hierarchy-growth} bounds
the approximation error at both ends uniformly, so the
convergence extends to \(\varphi\). On growing tori the
support-lift argument applies eventually to each compact
tuple. This suffices for the bound on the limit. At a finite
torus distant Euclidean supports may overlap after projection,
so a global Euclidean cutoff functional is not used there.

For a fixed torus take smooth cutoffs in a product of separated
coordinate neighborhoods. Expand a joint compact test in
torus Fourier modes and localize each slot by its cutoff.
At fixed \(n\) the per-support moment bound costs at most
\(C(1+|k|)^n\) on a mode of combined frequency
\(k\in\mathbb Z^{4n}\). Integrating by parts \(5n+1\) times
in a coordinate of maximal frequency gives a summable
bound after this loss, since the remaining exponent is
\(4n+1\). The series defines a local functional of finite
order. It agrees with the tuple moments by the same separated
tensor approximation as in the product-box proof. Agreement
on overlaps and a partition of unity give the unique
off-diagonal distribution.
\end{proof}


\section{Reflection forms and insertions}
\label{cl:sec-reflection}

Use one subsequence of the two-state moments
\eqref{cl:eq-cutoff-moments}. In this section \(G_n\) denotes its
limit with the given reference subtraction. The test spaces and
Schwartz seminorms are those in Definition~\ref{cl:def-flat-spaces}.
For an affine hyperplane \(\Pi\) let \(\theta_\Pi\) be reflection in
\(\Pi\), and let \(\Theta_\Pi\) denote simultaneous pullback on every
test variable. Choose a positive side of \(\Pi\).
The space \(\mathcal S_{+,n}(\Pi)\) consists of diagonal-flat
Schwartz tests also flat outside the set in which all variables
lie in that open side. Its degree-zero term is \(\mathbb C\).
We write \(\tau_v f(x)=f(x-v)\), with the simultaneous action on
multivariable tests. For \(\Pi=\{x_4=0\}\) the positive side is
\(x_4>0\).

\begin{proposition}[Reflection positivity of the limiting moments]
\label{cl:prop-reflection-positivity}
Suppose the two-state moments converge on the
separated tuples as in
Proposition~\ref{cl:prop-separated-compactness}. Let their
diagonal-flat extensions be those of
Proposition~\ref{cl:prop-flat-extension}. For every affine
coordinate or diagonal mirror \(\Pi\), and every finite sequence
\(F=(F_0,\ldots,F_N)\) with \(F_n\in\mathcal S_{+,n}(\Pi)\),
\begin{equation}
\label{cl:eq-reflection-form}
 q_{G,\Pi}(F,F)
 :=\sum_{m,n=0}^N
    G_{m+n}(\Theta_\Pi\overline{F_m}\otimes F_n)\ge0 .
\end{equation}
The same formula with a second sequence in the second slot
defines a Hermitian positive semidefinite form, antilinear in
the first sequence. It is continuous at each pair of fixed
finite degree bounds in the induced Schwartz topologies.
The completion of its quotient by the null space is separable.
\end{proposition}

\begin{proof}
First take finite linear combinations of compact smooth product
tests. In each product the factors have mutually separated
supports on the positive side. There is no separation
requirement between supports belonging to different products.
Choose lattice mirrors \(\Pi_j\) of the same normal type with
heights tending to that of \(\Pi\). Coordinate site and link
planes and diagonal planes are supplied by \cite[Theorem B]{LP}.
Every contributing plaquette in a positive product is wholly on
the positive side of \(\Pi_j\) for large \(j\), since its
diameter is \(O(a_j)\) and the support has a fixed positive
distance to \(\Pi\). Thus the corresponding polynomial in
\(X_j^\nu\) is a bounded function of the positive-side bonds at
the fixed cutoff.

The barycentre of a reflected plaquette is the reflected
barycentre. Its trace is unchanged by a change of starting
vertex or orientation. Consequently the raw observable
intertwines the lattice reflection and pullback of the test.
Invariance of \(\nu_j\) under \(\Pi_j\) gives the same identity
for the subtracted scalar. Apply lattice reflection positivity
in \(\rho_j\) to the entire polynomial, including its constant
term. Each entry of the resulting form is a moment of a
separated tuple: the left entries lie on the negative side,
the right entries on the positive side, and each side is
internally separated. Pullback by \(\theta_{\Pi_j}\) converges
to pullback by \(\theta_\Pi\) in \(C^1\) on a fixed compact
neighbourhood of each reflected support. The moving-test
conclusion of Proposition~\ref{cl:prop-separated-compactness}
therefore passes the finite sum to \eqref{cl:eq-reflection-form}.
Its degree-zero entry is \(G_0=1\).

For arbitrary \(F_m,F_n\) in the stated spaces, their mixed
tensor is diagonal-flat. At a coincidence within one group
every differentiated term vanishes by flatness of that group.
At a coincidence across groups, the common point cannot lie
strictly on the respective open sides for both factors.
Every derivative of at least one factor is therefore zero at
that point. The same argument covers zero time and all
multiderivatives. Affine pullback, conjugation and tensor
product are continuous in Schwartz topology. For example,
the product of the two weights of order \(r\) bounds the joint
weight of order \(r\). Thus the continuous extensions of
\(G_{m+n}\) apply to these mixed tensors. Approximate each
degree by finite sums of the separated positive-side tensors of
Proposition~\ref{cl:prop-flat-extension}, after an isometry
taking \(\Pi\) to the coordinate plane. Passing to the limit
in the finitely many terms gives the claimed inequality.
Polarization gives Hermiticity and Cauchy--Schwarz.
Schwartz space is separable and metrizable, and each flat
subspace is a separable subspace. Each finite direct sum
therefore has a countable dense set. Continuity of the form
makes its image dense for the quotient norm. The union of
these sets over the finite degree bounds is countable and
dense in the Hilbert completion. This proves separability.
\end{proof}

\begin{corollary}[Reflection positivity of growing-torus limits]
\label{cl:cor-torus-reflection}
Use the direct growing-torus compactness of
Corollary~\ref{cl:cor-torus-compactness}, with the graded
bounds and diagonal-flat extension of
Proposition~\ref{cl:prop-flat-extension}.
Then \eqref{cl:eq-reflection-form} holds on every affine
coordinate mirror for the self-centered limit.
If the torus sides also lie above the thresholds of
Proposition~\ref{cl:prop-large-torus-recognition},
it holds on every affine diagonal mirror as well.
\end{corollary}

\begin{proof}
Begin with the compact separated positive-side polynomials
in Proposition~\ref{cl:prop-reflection-positivity}.
Round a coordinate mirror to lattice mirrors of one fixed
type. The antipodal planes have the same type, since the
number of lattice sites along each side is even. For
large \(j\) the compact positive supports and their
sampled plaquettes lie strictly within the torus half
between these two planes. The first plane converges to
the prescribed mirror and the antipodal plane leaves
every fixed support. Finite-torus reflection positivity
from \cite[Proposition 3.1]{LP} gives the nonnegative
form on the entire lattice polynomial. Invariance
makes torus self-centering commute with reflection.
Use the fixed common lift of the supports and
the moving-test convergence to pass each form entry.

For a diagonal plane through the origin use the
approximate lattice inequality
\eqref{cl:eq-recognition-rp}. For a fixed polynomial of
degree at most \(N\), the supremum norm of the centered
lattice polynomial is bounded by \(C a_j^{-4N}\) for
all sufficiently small \(a_j\). This uses its fixed
compact supports and the plaquette count. Thus
\(\epsilon_{K_j}\mu_{\rm torus}(|F_j|^2)\to0\).
The limiting form is nonnegative. Translation invariance
carries it to every affine mirror of that orientation.
In both cases the separated tensor density and continuity
argument from Proposition~\ref{cl:prop-reflection-positivity}
gives the full diagonal-flat forms.
\end{proof}

For a mirror \(\Pi\), let \(\mathcal H_{G,\Pi}\) be the completion
of the finite-sequence quotient by the null space of
\(q_{G,\Pi}\). Write \([F]_G\) for a class when the mirror is
understood. The degree-zero class \(\Omega\) has norm one.
Compact separated positive-side tensor polynomials are dense
in this quotient. One can also use compact \(C^1\) factors:
smooth each within its fixed separated support neighbourhood,
and use the graded bound for the differences of the finitely
many norm entries.

\begin{proposition}[Translations on the reflection Hilbert space]
\label{cl:prop-reflection-semigroup}
Under Proposition~\ref{cl:prop-reflection-positivity}, fix one
of its mirrors and its inward unit normal \(e\).
Normal translations by \(ue\), \(u\ge0\), descend to a
strongly continuous semigroup \(T_u=e^{-uH}\) of positive
self-adjoint contractions, with \(H\ge0\).
Tangential translations descend to a strongly continuous
unitary group \(U_y=e^{iy\cdot P}\), \(y\perp e\).
The generators \(H\) and \(P\) have a joint spectral measure.
Both translation actions fix \(\Omega\).
\end{proposition}

\begin{proof}
The identity \(\Theta_\Pi\tau_{ue}=\tau_{-ue}\Theta_\Pi\)
and joint translation invariance give
\[
 q_{G,\Pi}(F,\tau_{ue}J)=q_{G,\Pi}(\tau_{ue}F,J).
\]
We first obtain a contraction on compact separated
polynomials. At a fixed cutoff, use a corresponding lattice
normal shift \(v\) into the chosen half-space. The same
reflection identity and semidefinite Cauchy--Schwarz give
\[
 \|\tau_v A\|_{\Pi_j}^2
     \le \|A\|_{\Pi_j}\,\|\tau_{2v}A\|_{\Pi_j}.
\]
The cutoff norm is at most the supremum norm of the bounded
cylinder function, uniformly over its translates at this
cutoff. If the initial norm is zero, the inequality gives
a zero shifted norm. Otherwise iterating with \(2^k v\) and
letting \(k\to\infty\) gives contraction. This argument uses
the cutoff supremum bound only before taking the cutoff limit.

Choose \(v_j\) in the lattice normal direction converging to
\(ue\), with nonnegative normal component. For a diagonal
normal one can use integer multiples of
\(a_j(e_i\pm e_l)\). Pass the two squared norms to the limit
by separated moving-test convergence as in the preceding
proof. This proves contraction on the compact tensor core.
Continuity of the reflected form and the translation on
Schwartz tests extend it to finite sequences and preserve the
null space. It therefore extends to the completed quotient.
The reflection identity gives self-adjointness, and
\[
 q_{G,\Pi}(F,\tau_{ue}F)=
 q_{G,\Pi}(\tau_{(u/2)e}F,\tau_{(u/2)e}F)
\]
gives positivity. Tangential translations and their inverses
preserve the form, so they descend to unitaries.
They commute with the normal translations by the translation law.

As the translation tends to zero, its tests converge in
Schwartz topology. Expand the squared norm of the difference
of two translated classes into form entries. Their continuity
gives strong continuity on the test classes, and the norm
bounds extend it to the completion. The positive contraction
semigroup has a nonnegative self-adjoint generator. In
particular \(T_1\) is injective. A vector in its kernel lies
in the kernel of every \(T_{2^{-k}}\) by self-adjointness
and the semigroup law, and strong continuity then gives
the zero vector. One may take \(H=-\log T_1\) by spectral
calculus, using positive roots and continuity to recover all
\(T_u\). The tangential translations commute with \(T_1\)
and hence with its spectral projections. The joint spectral
theorem for the strongly continuous abelian group
\((s,y)\mapsto e^{isH}U_y\) gives the asserted spectral
measure. The degree-zero class is fixed because it has no
translated arguments.
\end{proof}

For the coordinate time plane write \(q_G=q_{G,\Pi}\) and
\(\|F\|_G=q_G(F,F)^{1/2}\). For a compact test \(h\),
\(\phi(h)\) in a polynomial denotes the entry with test \(h\).
Write \(G(P)\) for the sum of moments of a finite polynomial
\(P\), including its constant term. In the expressions below
each moment has separated arguments unless the stated map
has already been extended by continuity. Use the source norm
in \eqref{cl:eq-sharp}.

\begin{proposition}[Uniform insertion on the actual reflection form]
\label{cl:prop-limit-insertion}
Use the same cutoff states, reference subtraction and
separated convergence as in
Proposition~\ref{cl:prop-reflection-positivity}. Suppose the
ordinary source and zero-source floor bounds and the
fixed-layer complete rooted read hypotheses of
\cite[Theorem C]{SI} hold uniformly for those data.
Let \(F,J\) be compact separated \(C^1\) tensor polynomials
with supports in positive time. For \(l\le1\) a power of
\(L\), and real \(h\in C_c^1\) supported in the radius-\(l/4\)
ball centered at spatial point \(y\in\mathbb R^3\) and time zero, one has
\begin{equation}
\label{cl:eq-insertion}
 \left|G\big(\Theta\overline J\,\phi(\tau_{(l/2)e_4}h)
                                  \tau_{le_4}F\big)\right|
 \le C\|h\|_\sharp\|J\|_G\|F\|_G .
\end{equation}
The operation \(\phi(\tau_{(l/2)e_4}h)\tau_{le_4}\)
descends to a bounded map of norm at most \(C\|h\|_\sharp\)
on the quotient completion. If \(0<\eta\le1\), \(0<c<1/64\),
and \(k\in C_c^1\) is supported in the radius-\(c\eta\) ball
about \(z\) with \(|z_4|\le\eta/8\), then for \(F,J\) strictly
above time \(\eta/4\) one has
\begin{equation}
\label{cl:eq-insertion-near}
 \begin{aligned}
 \left|G\big(\Theta\overline J\,\phi(k)F\big)\right|
 &\le C\eta^{-4}\|k\|_\sharp\,
     \|\tau_{-(\eta/4)e_4}J\|_G
     \|\tau_{-(\eta/4)e_4}F\|_G .
 \end{aligned}
\end{equation}
The constants are independent of the degree and number of
monomials of the surrounding polynomials, the states, the
cutoff sequence and the support positions. The assertions
also hold for the self-centered hierarchy \(G^0\) of
the same selected state sequence.
\end{proposition}

\begin{proof}
Use \cite[Proposition 11.2]{SI}, with its actual states
\(\rho_j,\nu_j\) and uniform source hypotheses. In its
lattice insertion the plane at \(a_j/2\) approximates time
zero, and the central cell has time coordinate \(a_j/2+l/2\).
Since \(l/a_j\) is odd for large \(j\), the center is a
lattice site. Choose its spatial center within \(O(a_j)\) of
\(y\). Translate all tests by this common small displacement
with time component \(a_j/2\), before translating the insertion
to the middle of its cell. The norm terms agree with the
lattice reflections of these shifted arguments, and the
passed fixed gaps give exactly the same limiting values as
in Proposition~\ref{cl:prop-reflection-positivity}.
The two-state estimate in the cited source uses the chosen \(\nu_j\) in the
inserted scalar as well. This verifies that
\eqref{cl:eq-insertion} applies on the present hierarchy
without a new centering or a new subsequence.

The polynomial \(\phi(\tau_{(l/2)e_4}h)\tau_{le_4}F\)
has strictly positive, internally separated arguments.
Use it as \(J\) in \eqref{cl:eq-insertion}.
When its norm is nonzero, divide the squared-norm bound by
that norm. This gives the bounded map and sends every null
class to zero. When its norm is zero the same norm bound
holds without division. For complex tests the estimate follows
by splitting real and imaginary parts and adjusting \(C\).

The near-mirror bound is \cite[Corollary 11.3]{SI}, which
uses the same limiting forms and the contraction of positive
normal translations. Its hypotheses keep the reflected left
and right entries at distance at least \(\eta/8-c\eta\)
from the insertion support. For clarity, the intermediate
slab estimate in that corollary gives
\[
 \|\phi(f)\tau_{4\eta e_4}F\|_G
 \le C(1+R)^3\eta^{-4}\|f\|_{C^1}\|F\|_G
\]
for tests \(f\) supported in \(\eta\le x_4\le3\eta\)
of diameter at most \(R\). It is a norm estimate with an
explicit time separation. The near-mirror inequality concerns
a sesquilinear form with the stated margins. Neither requires
a value for coincident insertions.

For \(G^0\) apply the same insertion result with both cutoff states
equal to \(\rho_j\). The self-centered moments converge on
the same subsequence by
Proposition~\ref{cl:prop-separated-compactness}, and the
hypotheses and constants are uniform for these choices.
\end{proof}

\section{Phase means and nontriviality}
\label{cl:sec-mean}

\begin{proposition}[Changing the finite phase mean]
\label{cl:prop-change-of-mean}
Suppose the selected full reference-centered hierarchy has
\(G_1(f)=b\int f\) with \(b\in\mathbb R\), and
\(G_n^0\) is its self-centered hierarchy on separated tuples.
Use their continuous flat extensions from
Proposition~\ref{cl:prop-flat-extension}, and suppose the
reflection form of Proposition~\ref{cl:prop-reflection-positivity}
is positive for at least one of them.
For \(A\subset\{1,\ldots,n\}\) let \(I_A\) integrate the
variables indexed by \(A\), retaining the other slots in
order. Then on \(\mathcal S_{\ne,n}\)
\begin{equation}
\label{cl:eq-mean-shift}
 G_n^0(F_n)=
 \sum_{A\subset\{1,\ldots,n\}}(-b)^{|A|}
                       G_{n-|A|}(I_AF_n).
\end{equation}
The inverse formula has \(b\) in place of \(-b\) and
interchanges \(G\) with \(G^0\). Both have zeroth moment one.
If one hierarchy has a bound \(A_0B^n(n!)^p|F_n|_{kn}\)
with \(p\ge0\), so does the other, with \(k\) replaced by
\(\max(k,6)\) and the exponential base enlarged by
\(C|b|\). Their reflection Hilbert spaces are unitarily
intertwined by the shift of tests in
\eqref{cl:eq-mean-shift}, preserving the vacuum and
translation actions. Joint Euclidean invariance and the
separated cluster property pass between the two hierarchies
whenever established on either one.
\end{proposition}

\begin{proof}
On a separated tensor, subtract \(b\int f\) in each
entry and expand the product. This gives
\eqref{cl:eq-mean-shift} including \(A=\{1,\ldots,n\}\),
whose remaining test is a scalar. Lemma~\ref{cl:lem-flat-marginal}
preserves the flatness of the remaining variables and gives
\(|I_AF_n|_r\le C^{|A|}|F_n|_{r+6|A|}\).
Continuity and density then identify the expression with
the diagonal-flat extension. If \(k\ge6\), take
\(r=k(n-|A|)\). The order is at most \(kn\).
The factorial at degree \(n-|A|\) is at most \((n!)^p\).
Summing the binomial coefficients bounds the new
exponential base by \(B+C|b|\), after increasing the
prefactor to cover degree zero.

On finite sequences let \(C_b\) send each degree-\(n\)
test to the lower-degree terms \((-b)^{|A|}I_AF_n\)
and sum all terms of each resulting degree.
The two opposite shifts compose to the identity:
for a fixed set of removed slots the coefficient in the
composition is the binomial sum of \(b\) and \(-b\).
The maps preserve positive-side flatness of the retained
variables. They also preserve ordered positive-time tests,
because any violation of strict order in the retained
times violates it in the original tuple. The maps commute
with conjugation and affine
isometries, since \(b\) is real and the integrated
Lebesgue measure is invariant. On a product of two
test groups, integration over subsets splits into the
independent subsets of those groups. It follows that
\[
 q_{G^0,\Pi}(F,J)=q_{G,\Pi}(C_bF,C_bJ).
\]
The invertibility of \(C_b\) proves positivity of both forms
from positivity of either one. The map on quotient classes
is therefore an isometry with an inverse on the test domain. Its unitary extension fixes
degree zero and intertwines the translations.
The same subset identities commute with common Euclidean
motions. In a clustering expression, the limit for each
pair of remaining groups factors, including an empty group,
and the finite subset sums then factor. This proves the
invariance and clustering claims.
\end{proof}

\begin{corollary}[The closed physical shape law]
\label{cl:cor-shape-law}
Let \(G\) be a separated reference-centered limit along
one admitted sequence of integrating and subtracting Wilson
states. Assume the running, source and reference hypotheses of
\cite[Theorem~A and Corollaries~6.5 and~7.1]{FP}, with the
finite-reference route or its cofinal alternative, and the same
physical normalization as in \eqref{cl:eq-observable}.
For each fixed sufficiently small \(d>0\), require the comparison
to apply at all sufficiently late cutoffs of this sequence to
every fixed pair in the following class. The tests \(f,h\) are
nonzero, nonnegative and compact \(C^1\). Each support lies in a
ball of radius \(d\), their mutual distance is in \([d,6d]\), and
each \(v=f,h\) satisfies
\[
 \int v\ge\vartheta_*\|v\|_\infty d^4,\qquad
 \operatorname{Lip}v\le\|v\|_\infty/(\vartheta_*d),
 \qquad \vartheta_*=\frac1{12000}.
\]
Write \(G_2^T(f,h)=G_2(f,h)-G_1(f)G_1(h)\), and let \(c_*>0\)
be the uniform inverse-coupling logarithmic coefficient in the cited comparison.
Put
\[
 \mathcal K(f,h)=\iint f(x)h(y)|x-y|^{-8}\,dx\,dy,
 \qquad
 M_d(f,h)=\frac{9\,\mathcal K(f,h)}
                 {\pi^4 c_*^2[\log(1/d)]^2}.
\]
There is a nonnegative function \(\epsilon_{\rm run}(d)\to0\)
such that
\begin{equation}\label{cl:eq-shape-envelope}
 \frac12(1-\epsilon_{\rm run}(d))^4
 \ \le\ \frac{G_2^T(f,h)}{M_d(f,h)}
 \ \le\ \frac32(1+\epsilon_{\rm run}(d))^4 .
\end{equation}
In particular, with \(m_v=\int v\) and \(M_v=\|v\|_\infty\),
\begin{equation}\label{cl:eq-shape-size}
 \frac{c\vartheta_*^2M_fM_h}{[\log(1/d)]^2}
 \ \le\ G_2^T(f,h)\
 \le\frac{C\min\{M_fM_h,d^{-8}m_fm_h\}}{[\log(1/d)]^2}.
\end{equation}
The function, constants and small-distance threshold can be
chosen uniformly over phases realized on these data when the
comparison assumptions have the corresponding all-state
uniformity.
\end{corollary}

\begin{proof}
For the late admitted cutoffs at scale \(d\), choose
\(\epsilon_{\rm run}(d)<1\) as a uniform upper bound for
\[
 |\zeta_K(d)-1|,\qquad
 |c_*g_K(d)^2\log(1/d)-1|.
\]
Such an envelope tends to zero by the two running limits in
\cite{FP}. Dividing the cutoff law by \(M_d(f,h)>0\)
gives the product of the squares of these two running factors
and \(1+r_{K,\rho}(f,h)\). The latter lies in \([1/2,3/2]\).
The connected correlation is unchanged by deterministic
reference subtraction. At a fixed pair its limit exists because
the two-point and both one-point functions converge.
Pass the resulting closed inequalities to that limit to obtain
\eqref{cl:eq-shape-envelope}. The bounds require no convergence
of the individual running factors as the cutoff is removed at
fixed \(d\).

Shrink the distance threshold so that
\(\epsilon_{\rm run}(d)\le1/2\). All point-pair distances on
the two supports lie in \([d,10d]\).
Hence \((10d)^{-8}m_fm_h\le\mathcal K(f,h)\le d^{-8}m_fm_h\).
Use the mass lower bound in the shape class and
\(m_v\le (\pi^2/2)d^4M_v\) to obtain
\eqref{cl:eq-shape-size}.
\end{proof}


\begin{corollary}[A nonzero scalar excitation]
\label{cl:cor-nontrivial}
Suppose a selected phase is realized by the convergent
separated reference-centered moments of
Proposition~\ref{cl:prop-separated-compactness}, with the flat
extension and positive reflection form of
Propositions~\ref{cl:prop-flat-extension}
and~\ref{cl:prop-reflection-positivity}. Let its realizing
states be admitted periodic Gibbs states
at the same bare sequence and physical source normalization
as in the two-point comparison of \cite[Theorem A and
Corollaries 6.5 and 7.1]{FP}. Require that at each fixed
sufficiently small \(d\) the chosen comparison route applies
for all sufficiently late cutoffs at that \(d\), with common
constants. The selected limit can retain a real mean
\(G_1(f)=b\int f\). For the fixed-shape real unit-mass tests
\(f_d\) below, the connected covariance satisfies
\begin{equation}
\label{cl:eq-nontrivial}
 c\,\frac{d^{-8}}{[\log(1/d)]^2}
 \le G_2(\theta f_d,f_d)-G_1(\theta f_d)G_1(f_d)
 \le C\,\frac{d^{-8}}{[\log(1/d)]^2}.
\end{equation}
Here \(\theta\) reflects in \(x_4=0\) and
\(c,C>0\) are independent of sufficiently small \(d\).
The reflection Hilbert space has a nonzero vector
orthogonal to the vacuum already at degree one.
\end{corollary}

\begin{proof}
Let
\(\varphi(x)=(12/\pi^2)(1-|x|^2)_+^2\) and
\(f_d(x)=d^{-4}\varphi((x-2de_4)/d)\).
Radial integration gives unit mass. The profile has
amplitude \(12/\pi^2\) and Lipschitz constant at most
\(8/(3\sqrt3)\) times that amplitude. The rescaled
positive and negative supports have radius \(d\) and
gap \(2d\). Use \(2d\) as the comparison scale. The supporting
balls have radius at most \(2d\), and the mass and Lipschitz
bounds of that shape class hold with parameter \(1/12000\).
The amplitude product is \((12/\pi^2)^2d^{-8}\), and
\(\log(1/(2d))\) is comparable to \(\log(1/d)\) as \(d\)
decreases. Apply the closed bounds
\eqref{cl:eq-shape-size} from Corollary~\ref{cl:cor-shape-law}
at scale \(2d\). They give \eqref{cl:eq-nontrivial}.
That corollary takes the cutoff limit at each fixed pair,
so this application removes the cutoff before sending \(d\) to zero.

In the reflection space, the squared norm of
\([\phi(f_d)]_G-G_1(f_d)\Omega\) equals the middle term of
\eqref{cl:eq-nontrivial}. The vector is orthogonal to
\(\Omega\), because \(q_G(1,\phi(f_d))=G_1(f_d)\).
It has positive norm and belongs to the completed ordered
positive-time sector. To see the latter, approximate the
degree-one test in \(C^1\) by smooth tests in a fixed
compact neighbourhood strictly above time zero. Their
classes converge by the separated moment bound. In a scalar
reconstruction that matches these moments it excludes a field proportional
to the identity.
\end{proof}


\section{Translation-ergodic limits}
\label{cl:sec-phases}

\subsection{The phase law at one spacing}

Let \(\mathcal I_\mu\) denote the completed translation-invariant
\(\sigma\)-field of a translation-invariant probability measure
\(\mu\). An ergodic measure is one for which \(\mathcal I_\mu\)
is trivial. Write \(\mathcal E_{a,\beta}\) for the ergodic
members of \(\mathcal G_{a,\beta}\).

We use the following external lattice estimate
\cite[uniform plateau estimate, Appendix~A]{Ext}.
Its hypotheses are translation invariance and reflection
positivity on the closed half-lattice for every coordinate
site, coordinate link and diagonal lattice mirror. Let
\(A,B\) be bounded Borel observables of bonds with both ends
in \(a[-r,r]^4\), for an integer \(r\ge1\). For \(v\in
a\mathbb Z^4\) and \(n=|v|_\infty/a-2r\ge1\), it gives
\begin{equation}\label{cl:eq-plateau}
 \begin{split}
 &\left|\mu(\overline A\,\tau_v B)-
    \mu\!\left(\overline{E_\mu[A\mid\mathcal I_\mu]}\,
                     E_\mu[B\mid\mathcal I_\mu]\right)\right|\\
 &\quad\le C\bigl[(r/n)^2+(r/n)^3\bigr]\,
              N_{\mu,r}(A)^{1/2}N_{\mu,r}(B)^{1/2},\\
 &N_{\mu,r}(A)=\max_{\lambda,\gamma}
      \mu\!\left(\overline{\theta_\lambda\gamma A}\,
                     \tau_{2ra e_\lambda}\gamma A\right).
 \end{split}
\end{equation}
Here \(\theta_\lambda\) reflects in \(x_\lambda=0\),
\(\gamma\) ranges over the coordinate sign changes, and
\(\tau_v\) translates supports by \(v\).
The terms in the maximum are squared reflection norms
at \(x_\lambda=ra\). The constant is numerical.
The lattice theorem in \cite[Theorem~B]{LP} supplies
all the reflection hypotheses for every member of
\(\mathcal G_{a,\beta}\).

\begin{proposition}[Ergodic Gibbs states]\label{cl:prop-fixed-pure}
Fix \(a>0\) and \(\beta\ge0\). Assume the uniform lattice
plateau estimate \eqref{cl:eq-plateau} for every member of
\(\mathcal G_{a,\beta}\).
The set \(\mathcal E_{a,\beta}\) is compact. It equals the set of
extreme points of \(\mathcal G_{a,\beta}\). Each
\(\mu\in\mathcal G_{a,\beta}\) is the barycentre of a unique
probability law \(D_\mu\) on \(\mathcal E_{a,\beta}\).
For every bounded Borel \(F\), the evaluation
\(\rho(F)\) under \(D_\mu\) has the law of
\(E_\mu[F\mid\mathcal I_\mu]\). The identification also holds
jointly for any finite list of such functions.
\end{proposition}

\begin{proof}
Let \(\gamma_D\) be the Wilson conditional kernel on a finite
bond set \(D\). If \(h\) is a bounded invariant density, then
\(\mu(hF)=\mu(h\gamma_DF)\). In fact, approximate \(h\) in
\(L^1(\mu)\) by bounded continuous cylinder functions. Translate
each approximant until its bonds avoid \(D\).
Invariance of \(h\) and \(\mu\) preserves its approximation
error. The translated approximant can be taken outside the
conditional integral defining \(\gamma_D\). The claimed
identity follows in the \(L^1\) limit. In particular,
conditioning on any invariant event of positive probability
preserves the Gibbs equations.

On the compact metric configuration space, a conditional
probability \(\omega\mapsto\mu^\omega\) given
\(\mathcal I_\mu\) can be constructed by taking conditional
expectations on a countable uniformly dense rational vector
space of real continuous functions containing \(1\). On one
set of full measure they are linear, positive on the
nonnegative members, have absolute value at most \(\|f\|_\infty\)
on a test function \(f\), and take value \(1\) on the constant \(1\). Their continuous
extensions are positive by uniform approximation and the
constant function. The Riesz representation theorem gives
probability measures. Approximation gives measurability on
continuous functions, and the bounded monotone class theorem
gives the conditional identity on bounded Borel functions.
Use countable continuous cylinder determining families in
the conditional identities for \(\gamma_D\) and the lattice
translations. The invariant-density argument and translation
invariance show that \(\mu^\omega\in\mathcal G_{a,\beta}\)
almost surely.

For a real continuous local function \(A\), apply
\eqref{cl:eq-plateau} to \(\mu\) and to each \(\mu^\omega\).
Dominated convergence in the disintegration of
\(\mu(A\tau_v A)\) shows that
\[
 \int\!\left[
   \mu^\omega\bigl(E_{\mu^\omega}
            [A\mid\mathcal I_{\mu^\omega}]^2\bigr)
             -\mu^\omega(A)^2\right]\,d\mu(\omega)=0 .
\]
The two integrands are measurable because the first is the
limit of local correlations. The integrand in brackets is
nonnegative. A countable dense family of real continuous
local functions is dense also in \(L^2\) for each component.
The vanishing conditional variances on this family prove
ergodicity almost surely. Push the conditional law to
\(\mathcal E_{a,\beta}\) to obtain \(D_\mu\).

A nontrivial invariant event splits its measure into two
different invariant Gibbs states by the density argument.
Conversely a constituent of a nontrivial convex decomposition
has a bounded invariant Radon--Nikodym density with respect
to the barycentre. An ergodic barycentre makes that density
constant. This proves the extreme point characterization.

For each fixed continuous local \(A,B\), the right side of
\eqref{cl:eq-plateau} has a state-independent bound using
their supremum norms. Hence a weak limit of ergodic Gibbs
states factorizes these two correlations at large separation.
The plateau identity and the same \(L^2\) argument show that
the weak limit is ergodic. This proves compactness.
Finally, a law on \(\mathcal E_{a,\beta}\) with barycentre
\(\mu\) has all its mixed evaluation moments determined by
\(\mu\). To obtain the product of \(q\) local evaluations,
translate the corresponding observables far apart, taking
successive limits with the other translates fixed.
The uniform plateau bound permits each limit through the
law. Products of continuous local evaluations form an algebra
that separates the compact set \(\mathcal E_{a,\beta}\).
Stone--Weierstrass gives uniqueness of the law.
\end{proof}

\subsection{A common cutoff sequence for the continuum phases}

Let \(X\) be a separated tuple of real compact \(C^1\) tests,
or a finite linear combination of such tuples.
The empty tuple is the unit.
Write \(X_j\) for its lattice polynomial formed with
\(X_j^\nu(f)\) from \eqref{cl:eq-observable}. Products of
tuples below are used when the resulting supports are
pairwise disjoint. Assume that the reference states
\(\nu_j\) define a separated limit \(G\), so that
\(\nu_j(X_j)\to G(X)\) for every such \(X\). This is one
choice of the integrating sequence in the separated
compactness proposition.

\begin{theorem}[The phase law of a reference limit]
\label{cl:thm-phase-law}
Let the reference states \(\nu_j\in\mathcal G_{a_j,\beta_j}\)
converge on the reference-centered separated compact moments
to \(G\) along one cutoff sequence, and assume the all-state
bound \eqref{cl:eq-graded-input} at these data. Use the
lattice plateau estimate \eqref{cl:eq-plateau}.
There is a compact metrizable space \(\mathcal P\) with
a probability law \(\lambda\),
with \(\operatorname{supp}\lambda=\mathcal P\), whose points
are distinct separated hierarchies \(G_\xi\).
For each \(\xi\in\mathcal P\) there are states
\(\rho_{j,\xi}\in\mathcal E_{a_j,\beta_j}\) along the given
cutoff sequence such that
\[
       G_\xi(X)=\lim_j\rho_{j,\xi}(X_j)
\]
for every separated compact tuple. Each \(G_\xi\) satisfies
the same graded estimate as \(G\). The evaluations on
compact tuples are continuous in \(\xi\), and
\begin{equation}\label{cl:eq-phase-barycentre}
                  G(X)=\int_{\mathcal P}G_\xi(X)\,d\lambda(\xi).
\end{equation}
For compact tuple polynomials \(X,Y\) and all sufficiently
large \(|v|_\infty\),
\begin{equation}\label{cl:eq-pure-cluster}
 |G_\xi(X\,\tau_vY)-G_\xi(X)G_\xi(Y)|
             \le C_{X,Y}|v|_\infty^{-2}.
\end{equation}
The constant and threshold are independent of \(\xi\).
The phase law is determined by \(G\) through the formula
\begin{equation}\label{cl:eq-phase-mixed}
 \int_{\mathcal P}\prod_{i=1}^q G_\xi(X^i)\,d\lambda(\xi)
 =\lim_{|v_q|\to\infty}\cdots\lim_{|v_2|\to\infty}
        G(X^1\,\tau_{v_2}X^2\cdots\tau_{v_q}X^q),
\end{equation}
where each inner limit has the outer translates fixed and
where the other, fixed groups are already separated.
Any compact Hausdorff representation with continuous
separating evaluations, a regular law of full support,
and these mixed moments is related to \((\mathcal P,\lambda)\)
by an evaluation-preserving measure-preserving homeomorphism.
\end{theorem}

\begin{proof}
For any fixed \(X,Y\), put their support sets in a physical
cube \([-R,R]^4\) with \(R\) a sufficiently large integer
and a positive margin to all its faces.
For \(a_j\le1\) put \(r_j=R/a_j\).
Each lattice polynomial reads only bonds in \(a_j[-r_j,r_j]^4\).
The reflected pair in \(N_{\rho,r_j}(X_j)\) is a linear
combination of products of separated compact tests with
a fixed positive gap at the common face. Its norm is
bounded uniformly in \(j,\rho\) by
\eqref{cl:eq-graded-input}. The same holds for \(Y\).
Sign changes, mirrors and lattice translations carry
\(X_j^\nu(f)\) to the corresponding transformed test since
\(\nu_j\) is invariant under these transformations.
For an ergodic \(\rho\), \eqref{cl:eq-plateau} now gives
\begin{equation}\label{cl:eq-lattice-pure-cluster}
 |\rho(X_j\,\tau_vY_j)-\rho(X_j)\rho(Y_j)|
                  \le C_{X,Y}|v|_\infty^{-2}
 \qquad (v\in a_j\mathbb Z^4,\ |v|_\infty\ge4R+1).
\end{equation}
Complex polynomials follow by using \(\overline X\) in the
plateau bound or by linearity.

Choose a countable real determining family of tuples. One
can use the compact supports of finite unions of rational
closed boxes with interiors, and a countable dense set of
\(C^1\) tests supported in each such union. Each separated
tuple is approximated in \(C^1\) in fixed separated unions.
This follows by first covering its compact supports in
separated open neighbourhoods by finite box unions.
The graded bound and multilinearity give equicontinuity
on these fixed containers, by telescoping in the entries.
Consequently convergence on the countable family implies
convergence on every separated tuple for these cutoff laws.

The evaluations on the determining family take values in
a countable product \(\mathcal K\) of compact real intervals,
by the graded bound. For \(\rho\in\mathcal E_{a_j,\beta_j}\)
let \(e_j(\rho)\in\mathcal K\) be its vector of evaluations.
Let \(\lambda_j=(e_j)_*D_{\nu_j}\). We show that the mixed
coordinate moments of \(\lambda_j\) are Cauchy. Given any
finite product of evaluations, first approximate the
product of its first two factors by the evaluation of a
single tuple with one group translated by a fixed large
integer vector, using \eqref{cl:eq-lattice-pure-cluster}.
Proceed with the third group and then each remaining
group. Each new translation is chosen after the existing
groups are fixed. All supports can therefore be kept
separated, all choices work uniformly in the state and
cutoff, and the total error can be made arbitrarily small.
The error terms are multiplied only by uniformly bounded
evaluations of the remaining factors. The resulting
single evaluation, integrated against \(D_{\nu_j}\),
converges by the assumed convergence of \(\nu_j\).
This proves the Cauchy assertion. Since the coordinate
algebra separates \(\mathcal K\), compactness and
Stone--Weierstrass give a unique weak limit \(\lambda\).

Put \(\mathcal P=\operatorname{supp}\lambda\). For
\(\xi\in\mathcal P\), every open neighbourhood of \(\xi\)
has positive \(\lambda_j\)-mass for all sufficiently large
\(j\). For a decreasing sequence of such neighbourhoods
of radius tending to zero, choose states in their inverse
images for all sufficiently large indices, increasing the
neighbourhood index along the full tail. This yields
\(e_j(\rho_{j,\xi})\to\xi\).
Uniform equicontinuity in the test functions extends the
limits to all separated tuples. It also proves continuity
of the resulting evaluations in \(\xi\), by uniform
approximation by the coordinate evaluations. The bound
and \eqref{cl:eq-phase-barycentre} follow by passage to
the limit, first on the countable family and then by
approximation.
The translation and signed coordinate symmetries pass on
each realizing sequence as in
Proposition~\ref{cl:prop-separated-compactness}.

Pass \eqref{cl:eq-lattice-pure-cluster} to the limit for
translations in \(\bigcup_{k\ge0}L^{-k}\mathbb Z^4\).
These are eventually lattice translations and form a dense
subgroup. Uniform continuity of the translated tests in
\(C^1\), with the separated containers enlarged slightly,
extends the estimate to every large real \(v\).
This proves \eqref{cl:eq-pure-cluster}.
To prove \eqref{cl:eq-phase-mixed}, apply
\eqref{cl:eq-phase-barycentre} and take the inner
translation to infinity. The other groups have compact
support and are kept fixed. Their product can be taken
as the fixed factor in \eqref{cl:eq-pure-cluster}.
The bound is uniform in \(\xi\). Repeat with the other
translations in the indicated order. For uniqueness,
embed both compact phase spaces in the product of real
evaluation coordinates indexed by all separated real
tuples. Each map is an embedding because its coordinates
are continuous and separate points. On the compact
union of their images, the two laws agree on the polynomial
coordinate algebra by \eqref{cl:eq-phase-mixed}.
Stone--Weierstrass gives equality of the regular laws.
Full support gives equality of the images, and the two
embeddings give the claimed homeomorphism.
\end{proof}

\begin{proposition}[The mean of a phase]\label{cl:prop-phase-mean}
There is a real continuous function \(b\) on \(\mathcal P\)
such that \(G_{\xi,1}(f)=b(\xi)\int f\) and
\(\int_{\mathcal P}b\,d\lambda=0\).
For each selected phase the hierarchies \(G_\xi\) and
\(G_\xi^0\), obtained from one another by adding or subtracting
the mean, are limits along the same realizing sequence.
More precisely \(G_\xi^0\) uses \(\rho_{j,\xi}\) for the
subtraction in \eqref{cl:eq-observable}, and
\begin{equation}\label{cl:eq-phase-mean}
 G_{\xi,n}(f_1,\ldots,f_n)
  =\sum_{I\subset\{1,\ldots,n\}}G^0_{\xi,|I|}(f_I)
            \prod_{i\notin I} b(\xi)\int f_i .
\end{equation}
For a fixed lattice plaquette \(p_j\) at each spacing,
\begin{equation}\label{cl:eq-phase-variance}
 \int b^2\,d\lambda
 =\lim_j 36a_j^{-8}
     \operatorname{Var}_{\nu_j}
         (E_{\nu_j}[A_{p_j}\mid\mathcal I_{\nu_j}]).
\end{equation}
In particular, for separated \(f,h\) the mixture two-point
function is the average of \(G^0_{\xi,2}(f,h)\) plus
\((\int b^2\,d\lambda)\int f\int h\).
\end{proposition}

\begin{proof}
Translations and signed coordinate permutations make
the plaquette mean the same at every plaquette of the
same spacing. Put
\(b_j(\rho)=6a_j^{-4}[\rho(A_{p_j})-\nu_j(A_{p_j})]\)
and \(R_j(f)=a_j^4\sum_p f(x_p)/6\).
Then \(\rho(X_j^\nu(f))=b_j(\rho)R_j(f)\).
The six shifted lattice Riemann sums give
\(R_j(f)\to\int f\). Use a real test \(f_*\) of integral
one in the determining family. Its evaluations and the
graded bound show that \(b_j\) is bounded uniformly for
large \(j\), and on a realizing sequence
\(b_j(\rho_{j,\xi})\to G_{\xi,1}(f_*)\).
This proves the singleton formula and continuity.
Equation~\ref{cl:eq-phase-barycentre} gives the zero
average. Expand the finite products after the change
of subtraction to obtain \eqref{cl:eq-phase-mean}.

The conditional identification in
Proposition~\ref{cl:prop-fixed-pure} gives
\(\int b_j(\rho)^2\,dD_{\nu_j}(\rho)\) equal to the
variance expression at spacing \(a_j\). Since
\(R_j(f_*)\to1\), weak convergence of the phase laws
on this bounded coordinate gives \eqref{cl:eq-phase-variance}.
The last assertion follows from \eqref{cl:eq-phase-mean}
at order two. The finite-spacing plateau formula also
identifies the variance before rescaling as the limit
of the connected plaquette correlation under far lattice
translation at that spacing.
\end{proof}


\section{Euclidean invariance of the limiting moments}
\label{cl:sec-euclidean}

\begin{proposition}[The joint transfer spectrum]
\label{cl:prop-cube-cone}
For a limit through the infinite-lattice states of
Proposition~\ref{cl:prop-reflection-positivity},
use the reflection Hilbert space of any coordinate
mirror. At each cutoff assume the external lattice carrier
of \cite[Uniform lattice plateau estimate and the Euclidean
Hilbert bridge]{Ext}: for joint one-step spectral values
\(e^{-\epsilon}\in(0,1]\) in the normal direction and
\(e^{ip}\), \(p\in[-\pi,\pi]\), in any tangential coordinate,
\(\sinh(\epsilon/2)\ge|\sin(p/2)|\).
Zero normal-transfer values are also permitted.
The nonnegative normal generator \(H\) and
the tangential generators \(P_1,P_2,P_3\) have joint
spectral support in
\begin{equation}\label{cl:eq-cube-cone}
                     E\ge\max_i |p_i|.
\end{equation}
\end{proposition}

\begin{proof}
Use the site mirror at zero and let \(T_j,V_{j,i}\) be
the normal transfer and the one-step tangential translations
at cutoff \(a_j\). The coordinate and diagonal reflection
hypotheses of the assumed carrier are supplied by
\cite[Theorem~B]{LP}. Apply it on the joint spectral
variables \(T_j=e^{-\epsilon}>0\) and \(V_{j,i}=e^{ip}\).
Convexity of \(\sinh\) on the nonnegative axis and
\(|\sin(kx)|\le k|\sin x|\) give, for every positive
integer \(k\),
\[
 (1-T_j^k)^2\ \ge\
          T_j^k(2-V_{j,i}^k-V_{j,i}^{-k}).
\]
The inequality includes the kernel of \(T_j\), where its
two sides are \(1\) and \(0\). Fix \(t>0\) and choose
\(k_j a_j\to t\). On a compact positive-time polynomial
each term is a reflected moment with at most a bounded
positive time shift and a bounded tangential shift.
The reflected groups retain their fixed separation.
The centered lattice observables intertwine these lattice
shifts and the site reflection. The graded bound and
moving-test convergence pass the inequality to
\[
 (1-e^{-tH})^2\ \ge\
       e^{-tH}(2-e^{itP_i}-e^{-itP_i})
\]
on the dense compact positive-time classes and hence on
the whole Hilbert space. Proposition~\ref{cl:prop-reflection-semigroup}
supplies the strongly commuting generators and \(H\ge0\).
For the countable sequence \(t=1/m\) the joint multiplier
\(2e^{-tE}[\cosh(tE)+\cos(tp_i)-2]\) is nonnegative
almost everywhere. Dividing the bracket by \(t^2\)
and taking \(m\to\infty\) proves \(E^2\ge p_i^2\)
on a common set of full spectral measure.
\end{proof}

\subsection{Analytic criteria for rotations}
\label{cl:subsec-rotation-criteria}

We use four external criteria from \cite[analytic insertion
densities, planar angular kernels, two-point angular elimination,
and all-orders angular closure]{Ext}. Their common input is a
centered scalar hierarchy on separated compact \(C^1\) tuples.
The hierarchy must be symmetric, real, translation invariant,
invariant under signed coordinate permutations, and continuous
on each separated support collection with position-uniform
bounds at fixed diameters and distances. Its coordinate
reflection forms must have the vacuum vector and the translation
actions of Proposition~\ref{cl:prop-reflection-semigroup}
with spectrum \eqref{cl:eq-cube-cone}. The criteria use the
following further data and implications.

\begin{enumerate}
\item The insertion-to-density result uses
\eqref{cl:eq-insertion-near} with a single exponent \(p\)
in place of \(4\), where \(p\ge0\), uniformly over all compact
outside polynomials. Fix \(0<c<1/64\) and \(0<\eta\le1\).
For any coordinate plane, write
\(y=(x_\lambda,x_\mu)\) with \(x_\lambda\) normal time
and \(z\) for the remaining coordinates. For a real compact
\(C^1\) transverse test \(\varphi(z)\), the insertion
\(\alpha(y)\varphi(z)\), with \(\alpha\) compact \(C^1\)
and supported in \(|x_\lambda|<c\eta\), has the bound
\[
 C_\varphi \eta^{-(p+5)}\|\alpha\|_{L^1}
         \|F\|_{(\eta/4)}\|J\|_{(\eta/4)}
\]
when the outside supports are strictly above \(\eta/4\).
Here the pulled-down norm is the reflection norm after
translation by \(-(\eta/4)e_\lambda\).
The result supplies compatible analytic operator densities
between positive time margins and translation identities
for their regularized point-plate operators.

\item The planar kernel criterion uses these \(L^1\) plate
bounds, the indicated spectral carrier in the plane, the
common local continuity and the signed coordinate
symmetries. It supplies holomorphic product kernels and
continued reflection Gram matrices. Their densities agree
with the compact separated moment functional on ordinary
planar ordering charts. In particular for all distinct
planar pairs there is a continuous two-point density
\(K_{\varphi,\psi}\) with
\[
 G_2^0(\alpha\otimes\varphi,\beta\otimes\psi)
  =\iint\alpha(y)\beta(y')K_{\varphi,\psi}(y'-y)\,dy\,dy'
\]
when the planar profiles have disjoint supports.
For \(t>2s>0\) this density in the coordinate frame
has the representation
\[
 K_{\varphi,\psi}(t,x)=
   \langle z_{\varphi,s},
     e^{-(t-2s)H+ixP}z_{\psi,s}\rangle,\qquad
 z_{\varphi,s'}=e^{-(s'-s)H}z_{\varphi,s}\quad(s'>s).
\]
At order \(n\) the angular coefficients are real analytic
away from planar coincidences and the band contains only
\(|m|\le Cn\) with characters \(e^{4imu}\) for planar
rotation angle \(u\). The continued Gram identities
give the highest-coefficient positivity on reflected
half-plane configurations and the inequality
\(2b_{k+l}\le b_{2k}+b_{2l}\) for the global largest
absolute angular indices \(b_j\). These statements
are conclusions of the kernel construction.

\item The order-two criterion uses the continuous
two-point density and the coordinate spectral
representation just stated. It also requires
positivity at the diagonal mirror whose normal makes
angle \(\pi/4\) with the time axis in that plane.
Its short-distance input is
\[
 0\le G^0_2(f,h)
      \le u(d)\|f\|_\infty\|h\|_\infty,\qquad u(d)\to0 ,
\]
uniformly on the shape class of
Corollary~\ref{cl:cor-shape-law}, at its fixed
parameter \(\vartheta_*=1/12000\), for every sufficiently
small \(d>0\).
It gives radiality for each equal-label planar kernel
and then for mixed labels by polarization. Thus \(b_2=0\).

\item The all-orders criterion uses the band, Gram
identities, top-coefficient positivity, the coordinate
reflection and cone hypotheses, and \(b_2=0\). It also uses
the off-diagonal distribution supplied by
Proposition~\ref{cl:prop-flat-extension}.
Its polygon and integer degree argument removes the
remaining harmonics. Its common-cut argument identifies
the invariant moment functional on arbitrary separated
compact four-dimensional tuples, including when planar
projections overlap. It yields invariance under each
coordinate-plane rotation, and hence under
\(\operatorname{O}(4)\) with the signed coordinate
symmetries.
\end{enumerate}

\begin{theorem}[Rotations of the realized limits]
\label{cl:thm-euclidean-rotations}
Let \(G\) be a separated limit through the stated
infinite-lattice Wilson states, with the uniform
ordinary-source and insertion assumptions of the
setting. Suppose its own realizing sequence satisfies
the assumptions of the two-point shape bound in
\cite[Theorem~A and the ultraviolet shape corollary]{FP},
using that theorem or its cofinal comparison
alternative at the same bare data. Use the external lattice
carrier of \cref{cl:prop-cube-cone} and the four analytic
criteria stated in \cref{cl:subsec-rotation-criteria}.
Then \(G^0\)
and \(G\) are invariant under the Euclidean group on
separated tuples and on \(\mathcal S_{\ne,n}\) at
each order. When the comparison assumptions are
uniform over all states at the fixed data, this
assertion applies to every phase in
Theorem~\ref{cl:thm-phase-law} and to its barycentre.
\end{theorem}

\begin{proof}
The self-centered hierarchy \(G^0\) is realized along
the same cutoff sequence as \(G\), with the integrating
states also used as the subtracting states.
Proposition~\ref{cl:prop-separated-compactness} and
\eqref{cl:eq-graded-input} give the scalar symmetries
and the required position-uniform \(C^1\) continuity.
Propositions~\ref{cl:prop-reflection-positivity} and
\ref{cl:prop-reflection-semigroup} give the coordinate
and diagonal positivity, the Hilbert translation actions
and the unit vacuum. Proposition~\ref{cl:prop-cube-cone}
gives the carrier in every coordinate picture.
Proposition~\ref{cl:prop-limit-insertion} on \(G^0\)
gives the insertion hypothesis with \(p=4\).
Consequently the first two criteria supply the plate
bound with \(p+5=9\), the planar density and its spectral
representation, and the all-orders band and Gram data.

Additive mean shifts leave connected two-point functions
unchanged. Corollary~\ref{cl:cor-shape-law} gives the
closed positive-shape law on the specified class with its
exact \(\vartheta_*\). Its upper bound can be taken to be
\(u(d)=C[\log(1/d)]^{-2}\) for small \(d\).
The sign follows from the lower bound on that class.
The diagonal mirror required by the order-two criterion
is one of the limiting diagonal lattice reflections.
Thus the third criterion gives \(b_2=0\), and the fourth
gives rotation invariance on all separated compact tuples.
Translations were already available. Restore \(b\int f\)
by the finite subset expansion, which commutes with every
isometry. The uniqueness and density of the flat Schwartz
extension in Proposition~\ref{cl:prop-flat-extension}
pass each invariance identity to the entire flat space.
In the phase-law application, every step uses the states
\(\rho_{j,\xi}\) from Theorem~\ref{cl:thm-phase-law}, or the
reference states themselves for the barycentre.
\end{proof}


\begin{corollary}[Isometries of the canonical phase law]
\label{cl:cor-phase-isometries}
Under Theorem~\ref{cl:thm-phase-law}, every Euclidean
isometry of the mixture \(G\) induces a measure-preserving
homeomorphism of \(\mathcal P\) acting by transport on
the evaluations. In particular, if \(G\) is
\(\operatorname{O}(4)\)-invariant then each \(G_\xi\)
is \(\operatorname{O}(4)\)-invariant.
\end{corollary}

\begin{proof}
Let \(r\) be an isometry of \(G\) and let \(R\) be
its linear part. Give the same phase space the transformed
evaluations \(X\mapsto G_\xi(r^{-1}X)\).
They are continuous and separate points.
Since \(r^{-1}\tau_v=\tau_{R^{-1}v}r^{-1}\),
the mixed moment formula \eqref{cl:eq-phase-mixed}
and invariance of \(G\) identify all their mixed moments
with those of the original evaluations. The uniqueness
assertion of Theorem~\ref{cl:thm-phase-law} gives a
measure-preserving homeomorphism carrying one system of
evaluations to the other. Separation of points gives the
group law for these homeomorphisms.

Each phase is invariant under the coordinate reflection
\(\theta_1\) by its lattice realization. If \(G\) is
rotation invariant, every rotate of \(G_\xi\) is another
phase by the first assertion, and hence is invariant under
\(\theta_1\). This says that \(G_\xi\) is invariant under
every conjugate of \(\theta_1\) by an orthogonal map.
These hyperplane reflections generate \(\operatorname{O}(4)\).
The same identities hold on the flat Schwartz extension
by density.
\end{proof}



\begin{proposition}[The zero-energy space and spatial rotations]
\label{cl:prop-phase-hilbert}
Use the realized phase law of Theorem~\ref{cl:thm-phase-law}
and the time reflection spaces of
Propositions~\ref{cl:prop-reflection-positivity}
and~\ref{cl:prop-reflection-semigroup}. Let \(H\) be the normal
generator in the reference mixture space and \(P_0\) its
projection onto \(\ker H\). On compact separated positive-time
polynomials the map
\[
 P_0[F]_G\longmapsto \bigl(\xi\longmapsto G_\xi(F)\bigr)
\]
extends to a unitary \(\ker H\to L^2(\mathcal P,\lambda)\)
sending \(\Omega\) to \(1\).
If \(\mathcal P\) has \(m<\infty\) points, then
\(\dim\ker H=m\). If \(\mathcal P\) is infinite, its zero-energy
space has countably infinite Hilbert dimension. In each
individual phase the kernel of the normal generator is the
vacuum line. Its spatial translation matrix elements tend
to the vacuum product for arbitrary reflection Hilbert
vectors.

For any of these hierarchies satisfying
Theorem~\ref{cl:thm-euclidean-rotations}, spatial
\(\operatorname{O}(3)\) acts by a strongly continuous unitary
representation on its reflection space. It fixes the vacuum,
commutes with the normal generator and rotates the spatial
translations. The joint spectral carrier in
\eqref{cl:eq-cube-cone} then improves to \(E\ge|p|\).
\end{proposition}

\begin{proof}
By spectral calculus \(T_t=e^{-tH}\) converges strongly
to \(P_0\) as \(t\to\infty\). For compact positive-time
polynomials \(F,J\), the barycentre formula and the
uniform phase cluster bound give
\[
 \begin{split}
 \langle[F]_G,T_t[J]_G\rangle
 &=\int_{\mathcal P}
   G_\xi(\Theta\overline F\,\tau_{te_4}J)\,d\lambda(\xi)\\
 &\longrightarrow\int_{\mathcal P}
             \overline{G_\xi(F)}\,G_\xi(J)\,d\lambda(\xi).
 \end{split}
\]
Each reflected tuple is separated. The identity
\(G_\xi(\Theta\overline F)=\overline{G_\xi(F)}\)
uses reality and reflection invariance of the phase.
Taking the spectral limit proves the asserted isometry
on the projection of the dense compact core.

The evaluation functions on that core contain constants,
are closed under conjugation and separate \(\mathcal P\).
For separation, translate any compact separating tuple
far enough into positive time and use phase translation
invariance. Products of evaluations are in their uniform
closure because
\[
 G_\xi(F)\,G_\xi(J)=
        \lim_{t\to\infty} G_\xi(F\,\tau_{te_4}J)
\]
uniformly in \(\xi\), by the same cluster bound.
For large \(t\) the product on the right has separated
positive-time arguments. The uniform closure is therefore
a conjugation-stable subalgebra of \(C(\mathcal P)\).
Stone--Weierstrass gives density in \(C(\mathcal P)\), and
hence in \(L^2(\mathcal P,\lambda)\). The isometry extends
onto this space. If \(\mathcal P\) is finite, each point has
positive mass by full support. If it is infinite, any finite
number of distinct points have disjoint open neighbourhoods
of positive measure. Their indicators give arbitrarily
large orthogonal sets. Since \(L^2\) of a compact metrizable
probability space is separable, its Hilbert dimension is
countably infinite.

For an individual phase the same normal cluster limit
gives the rank-one vacuum projection on the compact core,
and hence \(\ker H=\mathbb C\Omega\). With spatial
translations the compact cluster bound gives the same
matrix-element limit. Approximate each of two arbitrary
vectors by compact classes. Unitarity bounds the two
approximation errors independently of the spatial shift,
so the limit passes to the entire space.

For \(R\in\operatorname{O}(3)\), let the pullback act by
\(F((R^{-1}\mathbf x_1,t_1),\ldots,(R^{-1}\mathbf x_n,t_n))\).
Under Euclidean invariance this action preserves the positive-time
flat domain and commutes with time reflection. It is an
isometry for the reflection form, with inverse given by
the inverse rotation. Thus it descends to a unitary
\(V_R\). Spatial signed permutations already give these
unitaries under the coordinate symmetries alone. Pullback is continuous in the Schwartz topology
as \(R\) varies. The form then gives strong continuity on
finite test sequences, and the unitary norm gives it on
the completion. The exact actions on tests show
\(V_RT_t=T_tV_R\) and \(V_RU_yV_R^{-1}=U_{Ry}\), and
\(V_R\Omega=\Omega\). The joint spectral measure therefore
transforms by spatial rotation of its momentum variables.
The cube inequality \(E\ge |(Rp)_i|\) holds for each
rotation and coordinate. Intersect its supports over a
countable dense set of rotations, and use continuity in
\(R\) to obtain \(E\ge|p|\).
\end{proof}


The vacuum-line conclusion does not give a positive gap above
zero. For example, take \(H=0\oplus M_E\) on
\(\mathbb C\oplus L^2((0,\infty),Ee^{-E}\,dE)\), where
\(M_E\) is multiplication by \(E\). Its kernel is the first
summand, but the spectrum on the second reaches zero. The
vector \(v=0\oplus1\) has
\(\langle v,e^{-tH}v\rangle=\int_0^\infty Ee^{-(1+t)E}\,dE
=(1+t)^{-2}\). This polynomial decay in time coexists with
spectral support arbitrarily near zero.

\section{The corrected scalar reconstruction theorem}
\label{cl:sec-reconstruction}

For \(n\ge1\) write \(\mathcal S_{\mathrm{ord},n}\) for the
Schwartz functions all of whose derivatives vanish outside
\[
 \{(x_1,\ldots,x_n):0<(x_1)_4<\cdots<(x_n)_4\}.
\]
Set \(\mathcal S_{\mathrm{ord},0}=\mathbb C\). This space
is a subspace of \(\mathcal S_{+,n}\). The ordered
positive-time quotient is the Hilbert completion for the
restriction of the reflection form in
\eqref{cl:eq-reflection-form} to finite ordered sequences.

The growth condition below has content beyond the existence
of all moments. For a standard Gaussian \(Z\), the variable
\(X=e^Z\) has \(\mathbb E X^n=e^{n^2/2}\). For fixed
\(A,B,p>0\), the logarithm of \(AB^n(n!)^p\) grows only as
\(O(n\log n)\). Thus finiteness at every degree does not
supply the uniform factorial growth required for reconstruction.

\clneedspace{29\baselineskip}
\begin{proposition}[Applicability of scalar reconstruction]
\label{cl:prop-scalar-reconstruction}
Let the selected scalar hierarchy \(G_n\) have the
diagonal-flat extension of Proposition~\ref{cl:prop-flat-extension},
a positive full reflection form as in \eqref{cl:eq-reflection-form},
scalar permutation symmetry, and full Euclidean invariance.
Suppose its reflected compact separated tensor polynomials
cluster under all large spatial translations in nonzero
directions, with the product of the two means as their limit.
Assume there are \(A,B>0\), \(p\ge0\), and a nonnegative
integer \(k\), independent of \(n\), such that
\[
 |G_n(F)|\le A B^n(n!)^p |F|_{kn}
 \qquad(F\in\mathcal S_{\ne,n},\ n\ge1),
 \qquad G_0=1 .
\]
Then the corrected Osterwalder--Schrader theorem applies to
these scalar moments. It gives a Hermitian scalar Wightman
theory, unique up to unitary equivalence in its cyclic
realization on a separable Hilbert space, with the Wightman
covariance, positivity,
spectrum, locality and vacuum cluster properties and a
unique translation-invariant vacuum up to scalar multiples.
Its Euclidean Green functions agree with \(G_n\) on
\(\mathcal S_{\ne,n}\). If the two-point hypotheses of
Corollary~\ref{cl:cor-nontrivial} hold, its scalar field is
nontrivial. The construction applies with the reference
mean or with the self-centered hierarchy in
Proposition~\ref{cl:prop-change-of-mean}.
If \(G_1(f)=b\int f\) as in
Proposition~\ref{cl:prop-change-of-mean}, write \(\Phi\) for
the reference-centered reconstructed field. The self-centered
field on its cyclic polynomial domain is
\[
 \Phi^0(f)=\Phi(f)-b\left(\int_{\mathbb R^4} f(x)\,dx\right)I ,
 \qquad f\in\mathcal S(\mathbb R^4).
\]
The domain generated by the shifted field and the vacuum is
the same cyclic polynomial domain.
\end{proposition}

\begin{proof}
The corrected theorem is \(E'\to R'\) in
\cite[Section IV.1, pp.~287--288, equation (4.1)]{OS75}.
The Euclidean space in its linear growth condition is the
Schwartz subspace flat on every coincidence diagonal,
defined in \cite[Section II, p.~284]{OS75}. This is precisely
\(\mathcal S_{\ne,n}\). The radial weight and derivative cutoff in its
Section II, equation (2.1), use the convention of
\eqref{cl:eq-schwartz}. Choose an integer \(s\ge \max(1,k)\) and put \(B_+=\max(1,B)\).
Since \(n!\ge 2^{n-1}\) for \(n\ge1\), the sequence
\(A B_+ (n!)^{p+\lceil\log_2B_+\rceil}\) is a factorial
majorant and
\[
 |G_n(F)|
 \le A B_+(n!)^{p+\lceil\log_2B_+\rceil}|F|_{sn}.
\]
Thus the linear order and factorial coefficient are uniform
in the degree. If one instead uses the total-monomial
seminorm
\[
 q_{j,n}(F)=
 \max_{|\alpha|+|\gamma|\le j}
      \sup_{x\in\mathbb R^{4n}}|x^\gamma\partial^\alpha F(x)|,
\]
the same verification is quantitative. Expansion of
\((1+\sum_i|x_i|)^r\) bounds \(|F|_r\) by
\((1+4n)^r q_{2r,n}(F)\).
At \(r=sn\), use \(n^n\le e^n n!\) to bound the extra
factor by \(((5e)^s)^n(n!)^s\).
This again spends only a linear order and a factorial
coefficient. The growth hypothesis is \(E0'\).

For \(E1\) and \(E3\) use Euclidean invariance and scalar
permutation symmetry on the flat spaces. The positivity
test of \cite[Sections 2--3, pp.~86--88]{OS73} uses the
ordered positive-time sequences. Its reflection conjugates
and reverses the first group of variables. Permutation
symmetry removes the reversal inside each \(G_{m+n}\),
so the required form is the restriction of
\eqref{cl:eq-reflection-form}. This gives \(E2\).
The Hermiticity of the form, paired with degree zero, gives
\(G_n(\Theta\overline F)=\overline{G_n(F)}\) on positive-time
tests. Euclidean reflection invariance removes \(\Theta\).
Translate each compact separated tensor into positive time
and then use flat-test density. This proves the scalar
reality identity on the full diagonal-flat domain.
The ordered quotient embeds isometrically in
\(\mathcal H_{G,\{x_4=0\}}\). Its closed range contains
the vacuum. The spatial translations preserve strict time
order and act unitarily as shown below, so they preserve this
closed range. This comparison uses the ordered domain directly.

For \(E4\) take any two finite ordered Schwartz sequences.
Their classes in the full reflection space can be approximated
by compact separated positive-time tensor polynomials, by
Propositions~\ref{cl:prop-flat-extension}
and~\ref{cl:prop-reflection-positivity}.
Spatial translations preserve the reflection form by
translation invariance and have inverse translations.
They therefore act unitarily on the quotient and its
completion, as in Proposition~\ref{cl:prop-reflection-semigroup}. Replacing
either vector by its approximation changes any translated
matrix coefficient by at most the norm of the error times
the norm of the other vector, uniformly in the translation.
The same bound holds for the vacuum coefficients.
The compact reflected cluster law therefore passes to all
vectors of the full reflection space, in particular to the
ordered sequences. This extension is qualitative; it asserts
no common decay rate for arbitrary vectors. This is the spatially translated
cluster condition in \cite[Section 3, p.~88]{OS73}.
Thus the corrected theorem applies with \(E0'\) and
\(E1\)--\(E4\) to the same selected hierarchy.
In its cyclic representation the Wightman cluster
property gives the weak spatial-translation limit onto
the vacuum line, first on cyclic polynomial vectors and
then on all vectors by unitarity. A translation-invariant
vector lies in that line. The Wightman polynomial test maps
are continuous on each finite direct sum of Schwartz spaces.
Their countable dense subsets have dense image in the cyclic
Hilbert space, so that space is separable as well.

The scalar connected covariance in
\eqref{cl:eq-nontrivial} is unchanged by adding a constant
multiple of the identity to a field. A scalar field
proportional to the identity would give zero for it.
The mean shift of Proposition~\ref{cl:prop-change-of-mean}
commutes with the required symmetries and preserves
clustering and positivity, with the same form of growth.
This proves applicability for either centering.

On the cyclic domain, define \(\Phi^0\) by the displayed subtraction.
Expanding products with the identity shows that the cyclic
polynomial domains of \(\Phi\) and \(\Phi^0\) coincide.
The scalar integral is invariant under the Poincar\'e action
on tests, and the subtraction preserves Hermiticity and locality.
The vacuum and its translation representation are unchanged.
The same finite subset expansion transfers its cluster property.

For strictly time-ordered distinct Euclidean arguments,
the analytic correlations of the shifted field are obtained
by the finite subset expansion of its Wightman products.
Each retained subsequence stays in the ordered tube:
the imaginary difference between retained adjacent points
is a sum of the original forward differences. The constant
terms have the same Euclidean and Minkowski values.
Restriction to the Euclidean arguments therefore gives the
shift \eqref{cl:eq-mean-shift}. Every configuration of
distinct Euclidean arguments can be rotated to have distinct
time coordinates, by choosing an axis away from the finitely
many hyperplanes perpendicular to their pair differences.
Both Euclidean restrictions are symmetric and Euclidean
covariant. The equality in the ordered regions therefore
identifies the shift on compact tests away from the diagonals.
The continuous flat extensions identify it on \(\mathcal S_{\ne,n}\).
Uniqueness in the corrected reconstruction theorem now
identifies this shifted Wightman field with the reconstruction
from \(G^0\).
\end{proof}


\clearpage
\section{The scalar continuum conclusion}
\label{cl:sec-conclusion}

\begin{theorem}[A realized nontrivial scalar phase]
\label{cl:thm-scalar-continuum}
Let \(K_j\in\mathbb Z_{\ge0}\) and \(a_j=L^{-K_j}\downarrow0\)
give one sequence of admitted
real data with bare couplings \(\beta_j\ge0\), as in
\cref{cl:sec-setting}. Assume that the actual ordinary
source, rooted-domain and insertion hypotheses of
\cite[Theorems~B and C]{SI} hold uniformly on these data.
Assume also the state-uniform physical two-point hypotheses
of \cite[Theorem~A]{FP}, using its finite-reference route
or its cofinal alternative, with one common source
normalization and with the comparison applicable eventually
at each fixed sufficiently small physical separation.
Use the external plateau and spectral-carrier estimates
stated in \cref{cl:sec-phases,cl:prop-cube-cone} and the four
analytic criteria stated in \cref{cl:sec-euclidean}.

Choose reference states
\(\nu_j\in\mathcal G_{a_j,\beta_j}\), for example
fixed-spacing infinite-volume torus limits.
There is one subsequence with the following properties.
Its self-centered reference moments converge on every
separated compact \(C^1\) tuple to a hierarchy \(G\).
There is a compact metrizable full-support phase law
\((\mathcal P,\lambda)\) as in
\cref{cl:thm-phase-law}. Every \(\xi\in\mathcal P\)
is realized along this same subsequence by translation-ergodic
Wilson Gibbs states \(\rho_{j,\xi}\), with the original
reference subtraction \(\nu_j\).

The full phase moments \(G_{\xi,n}\) extend uniquely to
\(\mathcal S_{\ne,n}\). Their degree-zero value is one and,
with constants \(1\le A,B<\infty\) independent of \(n\) and \(\xi\),
\begin{equation}\label{cl:eq-main-growth}
 |G_{\xi,n}(\varphi)|
 \le A B^n(n!)^{k+32}|\varphi|_{20n},
 \qquad \varphi\in\mathcal S_{\ne,n}.
\end{equation}
Here \(k\) is the fixed exponent in
\eqref{cl:eq-graded-input}; the degree-zero convention is
that of \cref{cl:def-flat-spaces}.
The same estimate holds for the self-centered phase
hierarchies and the reference mixture after enlarging
\(A,B\). For every fixed flat test \(\varphi\), the
reference-centered evaluations \(\xi\mapsto G_{\xi,n}(\varphi)\)
are continuous, and their \(\lambda\)-integrals equal the
corresponding moments of the reference mixture.

Each phase is real, permutation symmetric and Euclidean
invariant on these test spaces. It is reflection positive
on every corresponding affine positive-time flat domain.
It clusters on separated compact groups, and its ordered
Schwartz sequences satisfy the spatial cluster condition.
Its first moment is \(b(\xi)\int f\), where \(b\) is a
bounded continuous real function with
\(\int b\,d\lambda=0\). The reference subtraction is
retained in this formula.

The full reflection Hilbert spaces have the zero-energy and
spatial rotation properties of \cref{cl:prop-phase-hilbert}.

The connected two-point functions satisfy the closed law
\eqref{cl:eq-shape-envelope} and the bounds
\eqref{cl:eq-shape-size} on the full shape class of
\cref{cl:cor-shape-law}, uniformly in \(\xi\).
For the fixed unit-mass positive-time profile \(f_d\) in
\cref{cl:cor-nontrivial}, the connected covariance in
every phase satisfies
\begin{equation}\label{cl:eq-main-signal}
 c\,\frac{d^{-8}}{[\log(1/d)]^2}
 \le
 G_{\xi,2}(\theta f_d,f_d)
       -G_{\xi,1}(\theta f_d)G_{\xi,1}(f_d)
 \le
 C\,\frac{d^{-8}}{[\log(1/d)]^2}
\end{equation}
for sufficiently small \(d>0\), with positive constants
and a threshold independent of \(\xi\).
The corrected scalar reconstruction gives each phase
a local Hermitian scalar Wightman field with positive
energy and a unique vacuum in its cyclic realization.
This field is nontrivial.
\end{theorem}

\begin{proof}
At each fixed datum \cref{cl:prop-fixed-spacing} gives
the required reference states. The actual source
hypotheses give \eqref{cl:eq-graded-input} for every
integrating/reference pair, by \cite[Theorem~B]{SI}.
Apply \cref{cl:prop-separated-compactness} to the pair
\((\nu_j,\nu_j)\). This chooses the single subsequence
used from now on. The reference singleton is zero.
The phase construction in \cref{cl:thm-phase-law}
uses the uniform bound before passing to the limit.
It supplies the realizing states on that same
subsequence and the uniform compact cluster law.

Apply \cref{cl:prop-flat-extension} with separation
exponent five to each realized hierarchy. The
constants depend on the common radius-one bound,
so they are uniform in the phase. Self-centering
uses the pairs \((\rho_{j,\xi},\rho_{j,\xi})\);
it therefore has the same graded estimate before
the limit. This proves \eqref{cl:eq-main-growth}.
Approximate a fixed flat test by finite sums of compact separated tensors
in the seminorm \(|\cdot|_{20n}\). The uniform bound makes
the phase evaluations converge uniformly in \(\xi\).
They are therefore continuous, and the compact barycentre
identity \eqref{cl:eq-phase-barycentre} passes to the flat test.
\Cref{cl:prop-phase-mean} identifies the mean and
its continuity and average.

The full reflected form follows from
\cref{cl:prop-reflection-positivity}. The insertion
estimate, spectral carrier and the physical shape
law meet the four analytic criteria, as checked in
\cref{cl:thm-euclidean-rotations}. Hence each phase
is Euclidean invariant. Conjugating a coordinate
reflection by an isometry transfers its positive
form to any affine mirror; the pullback preserves
the corresponding flat test domains.
The ordered Schwartz cluster condition follows
from the compact cluster law by the unitary
translation and Hilbert approximation argument
in \cref{cl:prop-scalar-reconstruction}.
The same compact cluster law and Euclidean invariance give
the Hilbert-space conclusions by \cref{cl:prop-phase-hilbert}.

For each fixed \(d\), the physical two-point bounds apply to all
late realizing states \(\rho_{j,\xi}\).
The fixed-pair limit argument of \cref{cl:cor-shape-law}
therefore gives the closed law on the full prescribed class.
Its specialization in \cref{cl:cor-nontrivial}
proves \eqref{cl:eq-main-signal}, with the common
constants inherited from the state-uniform bounds.
Finally \cref{cl:prop-scalar-reconstruction}
matches the quantitative growth, ordered
positivity, symmetry and cluster assumptions of
the corrected theorem. The positive connected
norm excludes a scalar field proportional to
the identity.
\end{proof}

The theorem is a deduction from the stated uniform
lattice estimates. It neither chooses contact terms
in the Euclidean moments nor identifies a complete
algebra of Yang--Mills observables. Such an
identification requires compatible joint continuum
limits for additional gauge-invariant observables.
Nontrivial scalar correlations alone also do not establish an
interaction criterion for an underlying field theory. If \(Z\)
is a standard Gaussian random variable, then \(Z^2-1\) has
cumulants \(2^{n-1}(n-1)!\) at every order \(n\ge2\), as follows
by differentiating \(-t-\tfrac12\log(1-2t)\) at zero. Thus a
quadratic observable of a Gaussian variable already has
nonzero higher connected moments.

\clneedspace{24\baselineskip}
\begin{thebibliography}{OS75}
\normalsize
\bibitem{LP} \emph{Periodic Wilson states and reflection positivity}, companion manuscript, 2026.
\bibitem{RT} \emph{Coupling trajectories and partition stability in lattice Yang--Mills theory}, companion manuscript, 2026.
\bibitem{SI} \emph{Scalar sources and uniform insertions for lattice Yang--Mills fields}, companion manuscript, 2026.
\bibitem{FP} \emph{Physical short-distance asymptotics of the lattice Yang--Mills action density}, companion manuscript, 2026.
\bibitem{Ext} \emph{Four-dimensional \(SU(2)\) Yang--Mills theory: A proof outline for \(\operatorname{tr}(F^2)\)}, 2026. General Euclidean, plateau, angular continuation and phase criteria, with the locators specified in the text.
\bibitem{Bal88} T.~Ba{\l}aban, \emph{Convergent renormalization expansions for lattice gauge theories}, Comm. Math. Phys. \textbf{119} (1988), 243--285. \href{https://doi.org/10.1007/BF01217741}{doi:10.1007/BF01217741}.
\bibitem{Bal89} T.~Ba{\l}aban, \emph{Large field renormalization. II. Localization, exponentiation, and bounds for the R operation}, Comm. Math. Phys. \textbf{122} (1989), 355--392. \href{https://doi.org/10.1007/BF01238433}{doi:10.1007/BF01238433}.
\bibitem{OS73} K.~Osterwalder and R.~Schrader, \emph{Axioms for Euclidean Green's functions}, Comm. Math. Phys. \textbf{31} (1973), 83--112. \href{https://doi.org/10.1007/BF01645738}{doi:10.1007/BF01645738}.
\bibitem{OS75} K.~Osterwalder and R.~Schrader, \emph{Axioms for Euclidean Green's functions II}, Comm. Math. Phys. \textbf{42} (1975), 281--305. \href{https://doi.org/10.1007/BF01608978}{doi:10.1007/BF01608978}.
\end{thebibliography}
\end{document}
